% Cayley's Collected Mathematical Papers, Vol. VII, Pages 251--300
% Transcribed from: collectedmathema07cayluoft.pdf

\documentclass[11pt,a4paper]{article}
\usepackage[utf8]{inputenc}
\usepackage[T1]{fontenc}
\usepackage{amsmath,amssymb,amsthm}
\usepackage{geometry}
\usepackage{titlesec}
\usepackage{fancyhdr}
\usepackage{enumitem}
\usepackage{array}
\usepackage{booktabs}
\usepackage{longtable}
\usepackage{microtype}
\tolerance=800
\emergencystretch=3em

\geometry{margin=1in}

\pagestyle{fancy}
\fancyhf{}
\fancyhead[C]{\small Cayley, Collected Mathematical Papers, Vol.\ VII}
\fancyfoot[C]{\thepage}
\renewcommand{\headrulewidth}{0pt}

\titleformat{\section}{\normalfont\large\bfseries}{\S\thesection.}{0.5em}{}
\titleformat{\subsection}{\normalfont\normalsize\itshape}{\thesubsection.}{0.5em}{}

\newcommand{\CayleyRef}[1]{[#1]}

\begin{document}

\begin{center}
  {\LARGE\bfseries Mathematical Papers}\\[0.5em]
  {\Large Arthur Cayley}\\[1em]
  {\large Volume VII, Pages 251--300}
\end{center}

\vspace{2em}


%% ========================================================================
%% PAPER 447 (continued): On the Rational Transformation Between Two Spaces
%% ========================================================================
\section*{\textbf{447.}}
\addcontentsline{toc}{section}{447. On the Rational Transformation Between Two Spaces}

[From the \textit{Proceedings of the London Mathematical Society}, vol.~iii (1869--1871), pp.~127--180. Read June 23, 1869.]

\textbf{\S 87.}\quad Considering next the cubic transformations, or those belonging to $n = 3$; in the case $m_1 = 4$, $h_1 = 4$, $\alpha_1 = 2$, the reciprocal transformation is of the order $9 - 4 = 5$; and in the case $m_1 = 5$, $h_1 = 5$, $\alpha_1 = 1$, the reciprocal transformation is of the order $9 - 5 = 4$: I do not consider these cases. But $m_1 = 6$, $h_1 = 7$, $\alpha_1 = 0$, the reciprocal transformation is of the order $9 - 6 = 3$; and assuming (as seems allowable) that the principal system does not contain any multiple point or curve, it must be of the same form as the original transformation, that is, we must have $n' = 3$, $m_1' = 6$, $h_1' = 7$, $\alpha_1' = 0$.

\textbf{\S 88.}\quad The transformations to be studied are thus,

$1^{\circ}$ The quadri-quadric transformation $n = 2$, $m_1 = 2$, $h_1 = 0$, $\alpha_1 = 1$, and $n' = 2$, $m_1' = 2$, $h_1' = 0$, $\alpha_1' = 1$; the principal system in each space consists of a point and of a conic (which may be a pair of intersecting lines); and the surfaces are quadrics.

$2^{\circ}$ The quadri-cubic transformation $n = 2$, $m_1 = 1$, $h_1 = 0$, $\alpha_1 = 3$, and $n' = 3$, $\alpha_1' = 0$, $m_1' = 3$, $h_1' = 3$, $m_2' = 1$, $h_2' = 0$: in the first space the principal system consists of three points and a line, and the surfaces are quadrics: in the second figure, the principal system consists of three simple lines and a double line and the surfaces are cubic surfaces passing through this principal system; that is, they are cubic scrolls.

$3^{\circ}$ The cubo-cubic transformation $n = 3$, $\alpha_1 = 0$, $m_1 = 6$, $h_1 = 7$, and $n' = 3$, $\alpha_1' = 0$, $m_1' = 6$, $h_1' = 7$; in each space the principal system is a sextic curve with seven apparent double points (but there are different cases to be considered according as the sextic curve does or does not break up into inferior curves), and the surfaces are cubic surfaces through the sextic curve.

\subsection*{The Quadri-quadric Transformation between Two Spaces.}

\textbf{\S 89.}\quad Starting from the equations $x' : y' : z' : w' = X : Y : Z : W$, we have here $X = 0$, \&c., quadric surfaces passing through a given point and a given conic (which may be a pair of intersecting lines). Take $x = 0$, $y = 0$, $z = 0$ for the coordinates of the given point; $w = 0$ for the equation of the plane of the conic; the conic is then given as the intersection of this plane by a cone having the given point for its vertex; or say the equations of the conic are $w = 0$, $(a, \ldots \mid x, y, z)^2 = 0$; the general equation of a quadric through the point and conic is
\[
w(ax + \beta y + \gamma z) + \lambda (a, \ldots \mid x, y, z)^2 = 0;
\]
and it hence appears that the equations of the transformation may be taken to be
\[
x' : y' : z' : w' = xw : yw : zw : (a, \ldots \mid x, y, z)^2;
\]
these give at once a reciprocal system of the same form; viz., the two sets are
\begin{align*}
x' : y' : z' : w' &= xw : yw : zw : (a, \ldots \mid x, y, z)^2, \\
x : y : z : w &= x'w' : y'w' : z'w' : (a, \ldots \mid x', y', z')^2.
\end{align*}

\textbf{\S 90.}\quad The Jacobian of the first space is at once found to be
\[
w^2(a, \ldots \mid x, y, z)^2 = 0;
\]
that of the second space is of course
\[
w'^2(a, \ldots \mid x', y', z')^2 = 0.
\]
The two spaces are similar to each other; we may say that there is in each of them a principal point and a principal conic; that the plane of the conic is the principal plane, and the cone having its vertex at the point and passing through the conic is the principal cone. To the principal point of either space corresponds any point whatever in the principal plane of the other space; and conversely. More definitely, the points of the one principal plane and the infinitesimal elements of direction through the principal point of the other space correspond according to the equations $x : y : z = x' : y' : z'$. To any point on the principal conic of either space corresponds in the other space, not a mere element of direction through the principal point of the other space, but a line of the principal cone; that is, to the points of the principal conic of the one space correspond the lines of the principal cone of the other space. The Jacobian of either space, consisting of the principal plane twice, and of the principal cone, is thus the principal counter-system of the other space.

\textbf{\S 91.}\quad (Writing $(a, \ldots \mid x, y, z)^2 = x^2 + y^2 + z^2$, $w = w' = 1$, the equations of transformation become
\[
x' : y' : z' : 1 = x : y : z : x^2 + y^2 + z^2,
\]
and
\[
x : y : z : 1 = x' : y' : z' : x'^2 + y'^2 + z'^2,
\]
or, what is the same thing, if for shortness
\[
x^2 + y^2 + z^2 = r^2, \quad x'^2 + y'^2 + z'^2 = r'^2,
\]
the equations are
\[
x' = \frac{x}{r^2},\quad y' = \frac{y}{r^2},\quad z' = \frac{z}{r^2},
\quad\text{and}\quad
x = \frac{x'}{r'^2},\quad y = \frac{y'}{r'^2},\quad z = \frac{z'}{r'^2},
\]
whence also $r'^2 = \frac{1}{r^2}$; this is the well known transformation by reciprocal radius vectors.)

\textbf{\S 92.}\quad The principal conic may be a pair of intersecting lines; taking its equations to be $w = 0$, $xy = 0$, the equations of transformation here become
\[
x' : y' : z' : w' = xw : yw : zw : xy,
\]
and
\[
x : y : z : w = x'w' : y'w' : z'w' : x'y'.
\]
There is no difficulty in the further development of the theory.

\subsection*{The Quadri-cubic Transformation between Two Spaces.}

\textbf{\S 93.}\quad It will be convenient to have the unaccented letters $(x, y, z, w)$ referring to the first space, and the accented letters $(x', y', z', w')$ to the second space. In the second space the equations $X' = 0$, $Y' = 0$, $Z' = 0$, $W' = 0$ are quadric surfaces passing through three fixed points (say the principal points) and through a fixed line (say the principal line) in the second figure. Taking $x' = 0$, $y' = 0$ for the planes passing through the principal line and through two of the principal points respectively; $z'$ for the plane passing through the three principal points, $w' = 0$ for an arbitrary plane passing through the first mentioned two principal points, the implicit factors of $x'$, $y'$, $w'$ may be so determined that for the third principal point $x' = y' = -w'$. That is, we shall have
\begin{align*}
&\text{for principal line}\quad x' = 0,\; y' = 0,\\
&\text{for principal points}\quad
\begin{cases}
(x' = 0,\; z' = 0,\; w' = 0),\\
(y' = 0,\; z' = 0,\; w' = 0),\\
(x' = y' = -w',\; z' = 0),
\end{cases}
\end{align*}
and this being so, the equation of a quadric surface through the principal points and line will be
\[
(\alpha x' + \beta y')z' + \gamma x'(y' + w') + \delta y'(x' + w'),
\]
and the equations of transformation may be taken to be
\[
x : y : z : w = x'z' : y'z' : x'(y' + w') : y'(x' + w').
\]

\textbf{\S 94.}\quad Writing these in the extended form
\begin{multline*}
x : y : z : w : x - y : z - w \\
= x'z' : y'z' : x'(y' + w') : y'(x' + w') : z'(x' - y') : w'(x' - y'),
\end{multline*}
and forming also the equation
\[
xy : (xw - yz) = z' : x' - y',
\]
we at once derive the reciprocal system of equations
\[
xw : yw : z : w = x(xw - yz) : y(xw - yz) : (x - y)xy : (z - w)xy,
\]
so that this is a cubic transformation. And the cubic surface in the first space (corresponding to an arbitrary plane $ax' + by' + cz' + dw' = 0$ of the second space) is
\[
(ax + by)(xw - yz) + c(x - y)xy + d(z - w)xy = 0;
\]
viz., this is a cubic surface having the fixed double line $(x = 0,\; y = 0)$, the fixed simple lines $(x = 0,\; z = 0)$, $(y = 0,\; w = 0)$, and $(x - y = 0,\; z - w = 0)$; it has also the variable simple line $(dz + cx = 0,\; dw + cy = 0)$. The principal figure of the first space thus consists of the three simple lines $(x = 0,\; z = 0)$, $(y = 0,\; w = 0)$, $(x - y = 0,\; z - w = 0)$, and of the line $(x = 0,\; y = 0)$, a double line counting four times in the intersection of two of the cubic surfaces.

\textbf{\S 95.}\quad The cubic surface as having the double line $(x = 0,\; y = 0)$ is a cubic scroll, the theory of which may be briefly explained. The scroll has a nodal directrix (viz., the double line), a simple directrix, and three fixed generating lines. Any plane through the nodal directrix meets the scroll in this line twice and in a conic; and the conic meets the nodal directrix in two points, and the simple directrix in one point. Taking an arbitrary plane and any conic therein, and taking for the simple directrix any line meeting the conic, and for the three generating lines the lines each meeting the nodal directrix and the simple directrix, and also the conic in one of its points respectively; the scroll is generated by a line which meets each of these five directrix lines. Or, what is the same thing, if through the nodal directrix and each of the three generating lines respectively we draw any plane, and in this plane a cubic passing through the intersection of the plane with the nodal directrix, and through the intersections of the plane with the three generating lines respectively; this cubic also passes through the intersection of the plane with the simple directrix. Conversely, if the plane be assumed at pleasure, and if, taking for the simple directrix any line which meets the given generating lines, we draw a cubic as above, then the scroll is the scroll generated by a line which meets each of the directrix lines, and also the cubic. If the plane be taken to pass through any generating line, then the cubic section breaks up into this line, and a conic; the conic does not meet the simple directrix, but it meets the nodal directrix; and any such conic will serve as a directrix; viz., the scroll is generated by the lines which meet the two directrix lines and the conic.

\textbf{\S 96.}\quad Any two scrolls as above meet in the three fixed generating lines, and in the nodal directrix counting four times; they consequently meet besides in a curve of the second order, which is a conic (one of the conies just referred to). In order to further explain the theory, suppose for a moment that the two scrolls had only a common nodal directrix; they would besides meet in a quintic curve; this curve would meet the nodal directrix in four points, viz., the points at which the two scrolls have a common tangent plane. Now if at any point of the nodal directrix the two scrolls have a common generating line, then the plane through this line and the nodal line is one of the two tangent planes of each scroll; that is, the scrolls have this plane for a common tangent plane. Hence, in the case of the common three generating lines, the points where these meet the nodal line are three of the four points just referred to; there remains therefore one point, which is the point where the conic meets the nodal line; through this point there are for each of the scrolls two generating lines; one of these for the first scroll, and one for the second scroll, lie in a plane through the nodal line, and are consequently the generating lines of the two scrolls respectively which pass through the remaining (one) intersection of the three scrolls.

\textbf{\S 97.}\quad Analytically we have the two equations
\begin{align*}
(ax + by)(xw - yz) + c(x - y)xy + d(z - w)xy &= 0, \\
(a_1x + b_1y)(xw - yz) + c_1(x - y)xy + d_1(z - w)xy &= 0,
\end{align*}
representing any two of the scrolls; and then putting for shortness $(ad_1 - a_1d) = \alpha$, $(bd_1 - b_1d) = \beta$, $(ac_1 - a_1c) = \gamma$, $(bc_1 - b_1c) = \delta$, the equations give
\[
(xw - yz)(\alpha x + \beta y) + (x - y)xy \cdot \gamma + (z - w)xy \cdot \delta = 0,
\]
which may be written
\[
(xw - yz)\{x(\alpha - \delta y) + y(\beta + \delta x)\} + \gamma(x - y)xy = 0,
\]
that is
\[
\alpha(xw - yz) + \beta yw + \gamma(x - y)y + \delta(xw - yz) = 0
\]
(the equation of the conic); and we thence obtain
\begin{align*}
xw - yz : (x - y)xy : xy &= \gamma(z - w) - \beta(x - y) : \delta(z - w) - \alpha(x - y) : \alpha(z - w) - \gamma w,\notag \\
&= X : Y : Z, \text{ say}.
\end{align*}
And the coordinates $x$, $y$, $z$, $w$ being here proportional to rational and integral functions of $X : Y : Z$, that is, to rational and integral functions of a single parameter, we see that the conic is a unicursal curve.

\textbf{\S 98.}\quad Any two scrolls as above meeting in a conic, a third scroll will meet the conic in six points; but these include the point on the nodal directrix twice, and the points on the three fixed generating lines each once; there is left a single point of intersection, viz., this is the one variable point of intersection of the three scrolls, which is in accordance with the theory.

\textbf{\S 99.}\quad For the Jacobian of the second space, we have
\[
\begin{vmatrix}
*, & z', & y' + w', & y' \\
z', & *, & x', & x' + w' \\
y' + w', & x', & *, & x' - y' \\
y', & x' + w', & x' - y', & *
\end{vmatrix} = 0;
\]
that is, $3x'y'z'(x' - y') = 0$; viz., $z' = 0$ is the plane containing the three principal points; and $x' = 0$, $y' = 0$, $x' - y' = 0$ are the planes which pass through the principal line and the three principal points respectively.

\textbf{\S 100.}\quad For the Jacobian of the first space, we have
\[
\begin{vmatrix}
2xw - yz, & -xz, & -xy, & x^2 \\
yw, & xw - 2yz, & -y^2, & xy \\
-xy, & x^2 - 2xy, & *, & (z - w)y \\
(z - w)y, & (z - w)x, & xy, & -xy
\end{vmatrix} = 0,
\]
that is, $3x^2y^2(x - y)^2(xw - yz) = 0$; viz., $x = 0$, $y = 0$, $x - y = 0$ are the planes through the nodal directrix and the three fixed generators respectively (each plane therefore occurring twice); and $xw - yz = 0$ is the quadric scroll generated by the lines which meet each of the three generators $(x = 0,\; z = 0)$, $(y = 0,\; w = 0)$, $(x - y = 0,\; z - w = 0)$; this scroll passing also through the nodal directrix $x = 0,\; y = 0$.

\subsection*{The Cubo-cubic Transformation between Two Spaces.}

\textbf{\S 101.}\quad The principal system in the first space is a sextic curve with 7 apparent double points; but this curve may be either a single proper sextic curve, or it may break up into inferior curves; and the like as regards the principal system in the second space. The cases which present themselves are:
\begin{enumerate}[label=(\Alph*)]
\item The principal system consists of six lines.
\item The principal system is a proper sextic curve; the lineo-linear transformation between two spaces.
\end{enumerate}

\textbf{(A) The Principal System consists of Six Lines.}

\textbf{\S 102.}\quad Taking in the first space, for the equations of the two lines, $(x = 0,\; y = 0)$ and $(z = 0,\; w = 0)$, and for the equations of the four lines, $(x = 0,\; z = 0)$, $(y = 0,\; w = 0)$, $(x - y = 0,\; z - w = 0)$, $(x - py = 0,\; z - qw = 0)$, then, if the equations of transformation are taken to be
\begin{multline*}
x' - py' : x' - y' : z' - qw' : z' - w' \\
= (x - py)(xw - yz) : (x - y)(qxw - pyz) : (z - qw)(xw - yz) : (z - w)(qxw - pyz);
\end{multline*}
these lead conversely (see post, No.~104) to a like system,
\begin{multline*}
x - py : x - y : z - qw : z - w \\
= (x' - py')M' : (x' - y')N' : (z' - qw')M' : (z' - w')N',
\end{multline*}
where for shortness
\begin{align*}
M' &= p(q - 1)x'w' - q(p - 1)y'z' + (pq - 1)(p - q)y'w', \\
N' &= (q - 1)x'w' - (p - 1)y'z' + (pq - 1)(p - q)y'w',
\end{align*}
or, as these are better written,
\begin{align*}
M' &= -q(p - 1)y'\{(p - 1)z' - (pq - 1)w'\} + p(q - 1)w'\{(q - 1)x' - (pq - 1)y'\}, \\
N' &= -(p - 1)y'\{(p - 1)z' - (pq - 1)w'\} + (q - 1)w'\{(q - 1)x' - (pq - 1)y'\}.
\end{align*}
Hence the principal system in the second plane is composed of the two non-intersecting lines $(x' = 0,\; y' = 0)$, $(z' = 0,\; w' = 0)$ and the four lines $\{(p - 1)z' - (pq - 1)w' = 0,\; (q - 1)x' - (pq - 1)y' = 0\}$, $(y' = 0,\; w' = 0)$, $(x' - y' = 0,\; z' - w' = 0)$, $(x' - py' = 0,\; z' - qw' = 0)$, each meeting each of the two lines.

\textbf{\S 103.}\quad The Jacobian of the first space is
\[
\begin{vmatrix}
2xw - yz - pyw, & -xz - pxw + 2pyz, & -y(x - py), & x(x - py) \\
2qxw - qyz - qyw, & -pxz - qxw + 2pyz, & -py(x - y), & x(x - y) \\
w(z - qw), & -z(z - qw), & xw - 2yz + qyw, & xz - 2qxw + qyz \\
qw(z - w), & -pz(z - w), & qxw - 2pyz + pyw, & qxz - 2qxw + pyz
\end{vmatrix} = 0,
\]
viz., this is $xyzw(x - y)(x - py)(z - w)(z - qw) = 0$, the equation of the planes each passing through one of the four lines.

\textbf{\S 104.}\quad To effect the foregoing transformation, writing
\begin{align*}
x' : y' : z' : w' &= (x - py)(xw - yz) : (x - y)(qxw - pyz) \\
&\qquad : (z - qw)(xw - yz) : (z - w)(qxw - pyz),
\end{align*}
or what will ultimately be the same thing, but it is more convenient for working with,
\begin{align*}
x' &= (x - py)(xw - yz), & y' &= (x - y)(qxw - pyz), \\
z' &= (z - qw)(xw - yz), & w' &= (z - w)(qxw - pyz),
\end{align*}
these give
\[
x - py = M'x',\quad x - y = N'y',\quad z - qw = M'z',\quad z - w = N'w',
\]
where $M'$, $N'$ are quantities which have to be determined; and thence
\begin{align*}
(1 - p)x &= M'x' - pN'y', & (1 - q)z &= M'z' - qN'w', \\
(1 - p)y &= M'x' - N'y', & (1 - q)w &= M'z' - N'w';
\end{align*}
whence also
\begin{multline*}
(1 - p)(1 - q)(xw - yz) = N'[\{(q - 1)x'w' - (p - 1)y'z'\}M' + (p - q)y'w'N'], \\
(1 - p)(1 - q)(qxw - pyz) = M'[-(p - q)x'z'M' + \{(pq - q)x'w' - (pq - p)y'z'\}N'];
\end{multline*}
but we have
\[
\frac{xw - yz}{qxw - pyz} = \frac{x'}{y'},\quad \frac{x - py}{x - y} = \frac{x'}{y'} \cdot \frac{M'x'}{N'y'} = \frac{M'}{N'};
\]
or, substituting,
\begin{multline*}
M'\{(q - 1)x'w' - (p - 1)y'z'\} + N'(p - q)y'w' \\
= M'\{-(p - q)x'z'\} + N'\{(pq - q)x'w' - (pq - p)y'z'\};
\end{multline*}
that is
\[
M'\{(q - 1)x'w' - (p - 1)y'z' + (p - q)x'z'\} = N'\{(pq - q)x'w' - (pq - p)y'z' - (p - q)y'w'\};
\]
or, what is the same thing,
\begin{align*}
M' &= (pq - q)x'w' - (pq - p)y'z' - (p - q)y'w', \\
N' &= (q - 1)x'w' - (p - 1)y'z' + (p - q)x'z';
\end{align*}
viz., $M'$, $N'$ having these values, the original equations
\begin{multline*}
x' : y' : z' : w' = (x - py)(xw - yz) : (x - y)(pxw - qyz) \\
: (z - qw)(xw - yz) : (z - w)(pxw - qyz),
\end{multline*}
give
\[
x - py : x - y : z - qw : z - w = M'x' : N'y' : M'z' : N'w'.
\]
If, in these equations, in place of $(x', y', z', w')$ we write $(x' - py',\; x' - y',\; z' - qw',\; z' - w')$, the new values of $M'$, $N'$ are found to be
\begin{align*}
M' &= p(q - 1)^2x'w' - q(p - 1)^2y'z' + (pq - 1)(p - q)y'w', \\
N' &= (q - 1)^2x'w' - (p - 1)^2y'z' + (pq - 1)(p - q)y'w',
\end{align*}
and we have the formulae of No.~102.

\textbf{(B) The Principal System of a Proper Sextic Curve; the Lineo-linear Transformation between Two Spaces.}

\textbf{\S 105.}\quad I start with the lineo-linear transformation, and show that this is in fact a transformation such that the principal system in either space is a sextic curve with seven apparent double points. I do not attempt any formal proof, but assume that the lineo-linear transformation is the most general one which gives rise to such a principal system.

We have between $(x, y, z, w)$, $(x', y', z', w')$ three lineo-linear equations; writing these first under the form
\begin{align*}
(P_1, Q_1, R_1, S_1 \mid x', y', z', w') &= 0, \\
(P_2, Q_2, R_2, S_2 \mid x', y', z', w') &= 0, \\
(P_3, Q_3, R_3, S_3 \mid x', y', z', w') &= 0,
\end{align*}
we have $x' : y' : z' : w' = X : Y : Z : W$, where $X$, $Y$, $Z$, $W$ are the determinants (each with its proper sign) formed out of the matrix
\[
\begin{matrix}
P_1, & Q_1, & R_1, & S_1 \\
P_2, & Q_2, & R_2, & S_2 \\
P_3, & Q_3, & R_3, & S_3
\end{matrix}
\]

\textbf{\S 106.}\quad Each of the surfaces $X = 0$, $Y = 0$, $Z = 0$, $W = 0$, or generally any surface
\[
aX + bY + cZ + dW = 0,
\]
is thus a cubic surface passing through the curve
\[
\begin{vmatrix}
-P_1, & Q_1, & R_1, & S_1 \\
P_2, & Q_2, & R_2, & S_2 \\
P_3, & Q_3, & R_3, & S_3
\end{vmatrix} = 0,
\]
which is at once seen to be of the order 6. In fact any two of these surfaces, for instance
\[
\begin{vmatrix}
P_1, & Q_1, & R_1 \\
P_2, & Q_2, & R_2 \\
-P_3, & Q_3, & R_3
\end{vmatrix} = 0
\quad\text{and}\quad
\begin{vmatrix}
P_1, & Q_1, & S_1 \\
P_2, & Q_2, & S_2 \\
-P_3, & Q_3, & S_3
\end{vmatrix} = 0,
\]
have in common a curve
\[
\begin{vmatrix}
P_1, & Q_1 \\
P_2, & Q_2 \\
P_3, & Q_3
\end{vmatrix} = 0,
\]
which is of the order 3; they consequently besides intersect in a curve of the order 6, which is the before mentioned curve of intersection of all the surfaces. And it further appears that the number of the apparent double points is $= 7$; in fact the formula in the case of two surfaces of the orders $\mu$, $\nu$, the complete intersection of which consists of a curve of the order $m$ with $h$ apparent double points, and of a curve of the order $m'$ with $h'$ apparent double points, the numbers $m$, $m'$, $h$, $h'$ are connected by the equation $2(h - h') = (m - m')(\mu - 1)(\nu - 1)$. (Salmon's \textit{Solid Geometry}, 2nd~Ed., p.~273 [Ed.~4, p.~311]). Hence, in the case of the two cubic surfaces intersecting as above (since for the cubic curve we have $m' = 3$, $h' = 1$, and for the sextic $m = 6$), the formula becomes $2(h - 1) = 12$, that is $h = 1 + 6 = 7$; or the number of apparent double points is $= 7$.

\textbf{\S 107.}\quad It thus appears that the principal system in the first plane is a curve of the order 6, with seven apparent double points: it is to be added that there are not in general any actual double points or stationary points, so that the class of the curve is $6 \cdot 5 - 2 \cdot 7 = 16$, and its deficiency is $\frac{1}{2} \cdot 5 \cdot 4 - 7 = 3$. For convenience I will refer to this as the curve $\Sigma$.

The transformation is obviously a symmetrical one; hence the principal system in the second space is a like sextic curve, which may be called $\Sigma'$.

\textbf{\S 108.}\quad Consider in the first space any point $P$ on the curve $\Sigma$; for this point the cubic surface passes through the line $L$, corresponding to the point $P'$ in the second space. It was seen, No.~42, that if the point $P$ is on $\Sigma$, then the line $L$ is an asymptote, or say it cuts $\Sigma$ twice; but I affirm further that the line $L$ cuts $\Sigma$ three times, or (what is the same thing) that the line $L$ is a line on the cubic surface. This being so, any plane through $L$ meets the cubic surface in $L$ and a conic; and the conic meets $\Sigma$ in $6$ points, but of these $3$ are on the line $L$, and there remain $3$ points, that is, any plane through $L$ meets $\Sigma$ in $3$ points; the locus of the lines $L$ is thus a scroll, say the scroll $\Pi$, and the curve $\Sigma$ is a triple line on the scroll.

\textbf{\S 109.}\quad The scroll $\Pi$ is the Jacobian of the first space; and as such it is of the order 8. (This is a theorem the analytical verification of which seems by no means easy). But without assuming the identity of the scroll $\Pi$ with this Jacobian, or taking the order of the scroll to be known, I proceed to show that the scroll $\Pi$ is the scroll generated by the lines each of which meets the curve $\Sigma$ three times; it will thereby appear that the order is $= 8$, and that the curve is a triple line on the scroll.

Consider a point $P'$ on $\Sigma'$, and the corresponding line $L$ of the first space: take $\theta'$ a plane in the second space; corresponding to it the cubic surface $\Theta$ in the first space. By imposing a single relation on the coefficients $(a, b, c, d)$ in the equation $ax' + by' + cz' + dw' = 0$ of the plane $\theta'$, we make it pass through the point $P'$; therefore by imposing this same single relation on the coefficients $(a, b, c, d)$ of the cubic surface $\Theta$, we make it pass through the line $L$, $\Theta$ is a cubic surface through $\Sigma$ and it is easy to see that the effect will be as above only if the line $L$ cuts the curve $\Sigma$ three times; this being so, the general cubic surface meets $L$ in three points (viz., the three intersections of $L$ with $\Sigma$), and if $\Theta$ be made to pass through a fourth point on the line $L$, it will pass through the line $L$; it thus appears that the line $L$ meets $\Sigma$ three times, and consequently that the scroll $\Pi$ is generated by the lines which meet $\Sigma$ three times.

\textbf{\S 110.}\quad The theory of a scroll so generated is considered in my ``Memoir on Skew Surfaces, otherwise Scrolls,'' [340] and [410]; viz., if $\Sigma$ is a given curve of the order $m$ and deficiency $q$, the scroll generated by the lines which meet $\Sigma$ three times is of the order $3m + 2q - 4$, having $\Sigma$ for a triple line; and the section by an arbitrary plane is a curve of that order having $3m + 2q - 4$ double points. In the present case $m = 6$, $q = 3$, and the order of the scroll is $= 18 + 6 - 4 = 20 - 4 = 8$; the section by a plane has 8 double points, and it is a unicursal curve of the order 8.

\textbf{\S 111.}\quad It may be remarked that any plane $\theta'$ meets the sextic curve $\Sigma'$ in six points; the corresponding six lines of the first space are generating lines of the scroll, lying in the plane $\Theta$ corresponding to the plane $\theta'$. Any point of $\Sigma'$ gives a line of the scroll; that is, the lines of the scroll correspond to the points of $\Sigma'$, or the scroll is related to $\Sigma'$ in the same manner that $\Sigma$ is to the scroll.

\textbf{\S 112.}\quad If the point $P'$ is on $\Sigma'$, then the corresponding line $L$ meets $\Sigma$ three times, and conversely; hence the scroll $\Pi$ (locus of the lines $L$) and the curve $\Sigma$ are related in the same manner as the curve $\Sigma'$ and the scroll $\Pi'$ (locus of the lines $L'$ corresponding to the points of $\Sigma$); viz., $\Sigma$ is a triple curve on $\Pi$, and $\Sigma'$ is a triple curve on $\Pi'$.

\textbf{\S 113.}\quad But to connect this with the Jacobian theory: the Jacobian of the first space is the locus of the points $P$ such that the first derivatives of $X$, $Y$, $Z$, $W$ (each of them a cubic function of $(x, y, z, w)$) all of them vanish; that is, it is the locus of the points such that the cubic surfaces through $\Sigma$ have a common tangent plane. If $P$ is on $\Sigma$, the cubic surfaces all of them pass through the line $L$, corresponding to the point $P'$ of $\Sigma'$; and if moreover they have a common tangent plane, then the line $L$ is on the Jacobian, that is, the scroll $\Pi$ generated by the lines $L$ lies on the Jacobian; and if the Jacobian is of the order 8, then the scroll $\Pi$ is the whole Jacobian.

\textbf{\S 114.}\quad To verify this, I consider the section by the plane $w = 0$. Writing $X$, $Y$, $Z$, $W$ for the cubic functions of $(x, y, z, w)$, and $(a, b, c, d)$ for arbitrary coefficients, the equation $aX + bY + cZ + dW = 0$ is that of any cubic surface through $\Sigma$; and if $(X, Y, Z, W)$ are each of them $= 0$, the cubic surface has the point $(x, y, z, w)$ for a double point. The section by $w = 0$ is a cubic curve $(a, b, c, d \mid X, Y, Z, W) = 0$, having the double points given by $X = 0$, $Y = 0$, $Z = 0$, $W = 0$, that is, the six intersections of $\Sigma$ by the plane $w = 0$; and the Jacobian is a surface of the order 8 having $\Sigma$ for a triple curve, or the section by $w = 0$ is a curve of the order 8, having the six points for triple points; that is, a curve of the order 8 with 6 triple points. The deficiency of a curve of the order 8 with 6 triple points is $= \frac{1}{2} \cdot 7 \cdot 6 - 6 \cdot 3 = 21 - 18 = 3$; the curve is thus unicursal, which agrees with the foregoing theory of the scroll $\Pi$.

\textbf{\S 115.}\quad Returning to the general theory, the equations of transformation are $x' : y' : z' : w' = X : Y : Z : W$; hence, reciprocally, for the transformation from the second space to the first, the equations are $x : y : z : w = X' : Y' : Z' : W'$, where $X'$, $Y'$, $Z'$, $W'$ are the determinants formed from the matrix of the coefficients of the lineo-linear equations taken in the reverse order; that is, from the matrix which expresses $(x', y', z', w')$ in terms of $(x, y, z, w)$. The whole is thus a perfectly symmetrical transformation between the two spaces; and we have in each space a sextic curve $\Sigma$ or $\Sigma'$ with 7 apparent double points, the principal system; and an octic scroll $\Pi$ or $\Pi'$, the Jacobian, having the sextic curve for a triple line.

\textbf{\S 116.}\quad The formula for the postulation of a multiple curve on a surface of the order $n$ is as follows: if the $r$-tuple curve consist of a simple curve of the order $m_1$ with $h_1$ apparent double points, of a double curve of order $m_2$ with $h_2$ apparent double points, and so on; and if moreover, the curves $m_1$, $m_2$ intersect in $k_{1,2}$ points, the curves $m_2$, $m_3$ in $k_{2,3}$ points, \&c.; then writing in general $\rho = \frac{1}{2}m(m - 1) - h$; that is, $\rho_1 = \frac{1}{2}m_1(m_1 - 1) - h_1$, $\rho_2 = \frac{1}{2}m_2(m_2 - 1) - h_2$, \&c., I find that the general condition of equivalence is
\begin{multline*}
\alpha_1 + (3n - 2)m_1 - 2\rho_1 \\
+ 8\alpha_2 + (12n - 16)m_2 - 16\rho_2 \\
+ \ldots + \alpha_r(3rn - 2r^2)m_r - 2r^2\rho_r = n^3 - 1;\\
- 5k_{1,2} - 8k_{2,3} - \ldots - (3r - 1)k_{r-1,r},
\end{multline*}
and that the general condition of postulation is
\begin{multline*}
\alpha_1 + (n + 1)m_1 - \rho_1 \\
+ \tfrac{1}{6}r(r + 1)(r + 2)\alpha_r \\
+ [\tfrac{1}{2}r(r + 1)n - \tfrac{1}{6}r(r + 1)(2r - 5)]m_r \\
- [\tfrac{1}{6}(r - 1)(r - 2)(r - 3)(r - 4) \\
\qquad\qquad + \tfrac{4}{3}r(r + 1)(2r + 1)]\rho_r \\
= \tfrac{1}{6}(n + 1)(n + 2)(n + 3) - 4;\\
- k_{1,2} - 2k_{2,3} - \ldots - \tfrac{1}{2}r(r - 1)k_{r-1,r}(s < r),
\end{multline*}
in which formulae it is however assumed that the curves have not any actual multiple points. This implies that if any one of the curves, say $m_r$, breaks up into two or more curves, the component curves are each to be simple curves, and the intersections of any two of them are to be actual double points, each counting as such, and not as apparent double points.

\textbf{\S 117.}\quad I will consider some particular cases of the formulae. First, let the $r$-tuple curve be a simple curve of the order $m_1$, with $h_1$ apparent double points; then $\rho_1 = \frac{1}{2}m_1(m_1 - 1) - h_1$, and the conditions are
\begin{align*}
\text{equivalence: } && \alpha_1 + (3n - 2)m_1 - 2\rho_1 &= n^3 - 1, \\
\text{postulation: } && \alpha_1 + (n + 1)m_1 - \rho_1 &= \tfrac{1}{6}(n + 1)(n + 2)(n + 3) - 4.
\end{align*}

\textbf{\S 118.}\quad Next let the $r$-tuple curve be a double curve of the order $m_2$, with $h_2$ apparent double points; then $\rho_2 = \frac{1}{2}m_2(m_2 - 1) - h_2$, and the conditions are
\begin{align*}
\text{equivalence: } && 8\alpha_2 + (12n - 16)m_2 - 16\rho_2 &= n^3 - 1, \\
\text{postulation: } && \alpha_2 + (n + 1)m_2 - \rho_2 &= \tfrac{1}{6}(n + 1)(n + 2)(n + 3) - 4.
\end{align*}

\textbf{\S 119.}\quad And so for a triple curve of the order $m_3$, with $h_3$ apparent double points; we have $\rho_3 = \frac{1}{2}m_3(m_3 - 1) - h_3$, and
\begin{align*}
\text{equivalence: } && 27\alpha_3 + (27n - 18)m_3 - 18\rho_3 &= n^3 - 1, \\
\text{postulation: } && \alpha_3 + (n + 1)m_3 - \rho_3 &= \tfrac{1}{6}(n + 1)(n + 2)(n + 3) - 4.
\end{align*}

\textbf{\S 120.}\quad It is to be observed that the postulation equations give at once the values of $\rho_1$, $\rho_2$, $\rho_3$, \&c.; but that the equivalence equations give also the values of $\alpha_1$, $\alpha_2$, $\alpha_3$, \&c. In particular, for a triple line ($m_3 = 1$, $h_3 = 0$, $\rho_3 = 0$), the postulation equation gives
\[
\alpha_3 = \tfrac{1}{6}(n + 1)(n + 2)(n + 3) - 4 - (n + 1) = \tfrac{1}{6}(n + 1)(n + 2)(n + 3) - (n + 5),
\]
and the equivalence equation gives
\[
27\alpha_3 + 27n - 18 = n^3 - 1, \quad\text{that is}\quad 27\alpha_3 = n^3 - 27n + 17.
\]

\textbf{\S 121.}\quad For a triple conic ($m_3 = 2$, $h_3 = 0$, $\rho_3 = 1$), we have
\begin{align*}
\text{postulation: } && \alpha_3 &= \tfrac{1}{6}(n + 1)(n + 2)(n + 3) - 4 - 2(n + 1) + 1 \\
&&&= \tfrac{1}{6}(n + 1)(n + 2)(n + 3) - 2n - 5, \\
\text{equivalence: } && 27\alpha_3 + 54n - 36 - 18 &= n^3 - 1, \\
&& 27\alpha_3 &= n^3 - 54n + 53.
\end{align*}

\textbf{\S 122.}\quad Let the triple curve be a pair of non-intersecting lines; here $m_3 = 2$, $h_3 = 0$, $\rho_3 = 1$, and the formulae are the same as for a triple conic; but we must consider the line-pair as a degenerate conic, and apply the formula with $\rho_3 = 1$. If the lines intersect, then $h_3 = 0$, but the intersection counts as an actual double point; the formula gives $\rho_3 = 1 - \frac{1}{2} = \frac{1}{2}$, which is fractional, and the case is thus excluded from the general formula.

\textbf{\S 123.}\quad Let the $r$-tuple curve consist of three right lines meeting in a point: this is an actual triple point, and the formulae do not apply. But calculating the postulation-terms by the formula, we have $m_r = 3$, $\rho_r = \frac{1}{2} \cdot 3 \cdot 2 - 0 = 3$; and the terms are
\[
[\tfrac{1}{2}r(r + 1)n - \tfrac{1}{2}r(r + 1)(2r - 5)] \cdot 3 - [\tfrac{1}{6}(r - 1)(r - 2)(r - 3)(r - 4) + \tfrac{4}{3}r(r + 1)(2r + 1)],
\]
which are
\[
= \tfrac{1}{2}r(r + 1)(3n - 4r + 4) - \tfrac{1}{6}(r - 1)(r - 2)(r - 3)(r - 4),
\]
or say
\[
= \tfrac{1}{2}r(r + 1)(3n - 4r + 4) + \tfrac{1}{6}(-r^4 + 10r^3 - 35r^2 + 50r - 24).
\]
I have found by an independent investigation that this value requires the correction
\[
+ \tfrac{1}{6}[r^4 - 8r^3 + 30r^2 - 56r + 24 + \tfrac{1}{2}\{1 - (-1)^r\}],
\]
and that the true value of the postulation is
\[
= \tfrac{1}{2}r(r + 1)(3n - 4r + 6) + \tfrac{1}{6}(2r^3 - 5r^2 - 6r + \tfrac{1}{2}[1 - (-1)^r]),
\]
viz., that this is the number of the conditions to be satisfied that a surface of the order $n$ may have for an $r$-tuple curve three given right lines meeting in a point.

\textbf{\S 124.}\quad The formula for the postulation of an $r$-tuple point is $\frac{1}{6}r(r + 1)(r + 2)$; and for the equivalence, $r^3$; but these formulae assume that there is no $r$-tuple curve through the point. If there be an $s$-tuple curve through the point ($s < r$), then the formulae require modification; and the investigation of the general formulae for the postulation and equivalence of multiple points and curves is a subject of considerable interest.

\clearpage

%% ========================================================================
%% PAPER 448: Note on the Cartesian with Two Imaginary Axial Foci
%% ========================================================================
\section*{\textbf{448.}}
\addcontentsline{toc}{section}{448. Note on the Cartesian with Two Imaginary Axial Foci}

\begin{center}
\textbf{NOTE ON THE CARTESIAN WITH TWO IMAGINARY AXIAL FOCI.}
\end{center}

[From the \textit{Proceedings of the London Mathematical Society}, vol.~iii (1869--1871), pp.~181, 182. Read June 9, 1870.]

Let $A$, $A'$, $B$, $B'$ be a pair of points and antipoints viz.;
\begin{itemize}
\item[] $(A, A')$ the two imaginary points, coordinates $(+\beta i, 0)$, $(-\beta i, 0)$;
\item[] $(B, B')$ the two real points, coordinates $(0, +\beta)$, $(0, -\beta)$;
\end{itemize}
and write $\rho$, $\rho'$, $\sigma$, $\sigma'$ for the distances of a point $(x, y)$ from the four points respectively:
\[
\rho = \sqrt{(x - i\beta)^2 + y^2},\quad \sigma = \sqrt{x^2 + (y + \beta)^2},
\]
\[
\rho' = \sqrt{(x + i\beta)^2 + y^2},\quad \sigma' = \sqrt{x^2 + (y - \beta)^2}.
\]

We have
\[
\rho^2 + \rho'^2 = 2x^2 + 2y^2 - 2\beta^2 = \sigma^2 + \sigma'^2 - 4\beta^2,
\]
\[
\rho\rho' = \sqrt{(x^2 + \beta^2 + y^2)^2 - 4\beta^2 y^2} = \sigma\sigma';
\]
and thence
\begin{align*}
(\rho + \rho')^2 &= (\sigma + \sigma')^2 - 4\beta^2, \\
(\rho - \rho')^2 &= (\sigma - \sigma')^2 - 4\beta^2;
\end{align*}
or say
\begin{align*}
\rho + \rho' &= \sqrt{(\sigma + \sigma')^2 - 4\beta^2}, \\
i(\rho - \rho') &= \sqrt{4\beta^2 - (\sigma - \sigma')^2}.
\end{align*}

The equation of a Cartesian having the two imaginary axial foci $A$, $A'$ is
\[
(p + qi)\rho + (p - qi)\rho' + 2\sqrt{q}c = 0;
\]
viz., this is
\[
p\sqrt{(\sigma + \sigma')^2 - 4\beta^2} + q\sqrt{4\beta^2 - (\sigma - \sigma')^2} + 2\sqrt{q}c = 0,
\]
or, what is the same thing, it is
\[
p\sqrt{(\sigma + \sigma')^2 + 4\beta^2} + q\sqrt{4\beta^2 - (\sigma - \sigma')^2} + 2c\sqrt{q} = 0,
\]
which is the equation expressed in terms of the distances $\sigma$, $\sigma'$ from the non-axial real foci $B$, $B'$. Of course, the radicals are to be taken with the signs $+$. This equation gives, however, the Cartesian in combination with an equal curve situate symmetrically therewith in regard to the axis of $y$.

The distances $\sigma$, $\sigma'$ may conveniently be expressed in terms of a single variable parameter $\theta$; in fact, we may write
\[
\pm\rho = \sqrt{(\sigma + \sigma')^2 - 4\beta^2} = k + k'\theta,
\]
that is
\[
4\rho^2 - (\sigma - \sigma')^2 = k'^2 - (k'\theta)^2 = k'^2(1 - \theta^2);
\]
And therefore
\[
\sigma + \sigma' = \pm\frac{1}{2}\sqrt{4\beta^2 + (k + k'\theta)^2},\quad
\sigma - \sigma' = \pm\frac{1}{2}\sqrt{4\beta^2 - (k - k'\theta)^2};
\]
so that, assigning to $\theta$ any given value, we have $\sigma$, $\sigma'$, and thence the position of the point on the curve. We may draw the hyperbola $y^2 = 4\beta^2 + c^2 x^2$, and the ellipse $y^2 = 4\beta^2 - c^2 k'^2 x^2$; and then measuring off in these two curves respectively the ordinates which belong to the abscissae $k + \theta$ for the hyperbola, $k - \theta$ for the ellipse, we have the values $\sigma + \sigma'$ and $\sigma - \sigma'$, which determine the point on the curve. Considering $k$, $p$, $q$, $c$ as disposable quantities, the conics may be any ellipse and hyperbola whatever, having a pair of vertices in common; and the complete construction is,---From the fixed point in the axis of $x$, measure off in opposite directions the equal distances $KM$, $KN$, and take
\begin{align*}
M &= \sigma + \sigma' \text{ the ordinate at in the hyperbola}, \\
N &= \sigma - \sigma' \quad\;\,\text{the ordinate at in the ellipse},
\end{align*}
where $\sigma$, $\sigma'$ denote the distances of the required point from the fixed points $B$ and $B'$ respectively, the distance of each of these from the origin being $= \beta$ the common semi-axis. We may imagine $N$ travelling from one extremity of the $x$-axis of the ellipse to the other, the value of $\sigma + \sigma'$ will be real and greater than $BB'$, that of $\sigma - \sigma'$ real and less than $BB'$, and the point $(\sigma, \sigma')$ will be real. The construction gives, it will be observed, the two symmetrically situated curves.

The $x$-semi-axis of the ellipse is $\frac{1}{c}\sqrt{2\beta}$, and the form of the curve depends chiefly on the value of the ratio $k : k'\sqrt{p}$; or, what is the same thing, $k'^2 : 2\beta q$. We see, for instance, that, in order that the curve may meet the axis of $y$ in two real points between the foci, the value $\theta = -k$ must give a real value of $\sigma - \sigma'$; viz., that we must have $4\beta^2 > 4k'^2$; that is, $\frac{1}{q}c^2 > k'^2$, or $k'^2 c^2 > q$. If $k$ has this value, viz., $k = \frac{1}{c}\sqrt{2\beta} = \frac{1}{c}\sqrt{\text{semi-axis}}$, the curve touches the axis of $y$ at the origin; if $k < \frac{1}{c}\sqrt{\text{semi-axis}}$, the curve cuts the axis of $y$ in two real points between the foci; if $k > \frac{1}{c}\sqrt{\text{semi-axis}}$, the curve does not cut the axis of $y$ between the foci.

\clearpage

%% ========================================================================
%% PAPER 449: Sketch of Recent Researches Upon Quartic and Quintic Surfaces
%% ========================================================================
\section*{\textbf{449.}}
\addcontentsline{toc}{section}{449. Sketch of Recent Researches Upon Quartic and Quintic Surfaces}

\begin{center}
\textbf{SKETCH OF RECENT RESEARCHES UPON QUARTIC AND QUINTIC SURFACES.}
\end{center}

[From the \textit{Proceedings of the London Mathematical Society}, vol.~III (1869--1871), pp.~186--195. Read Nov.~10, 1870.]

The classification of quartic surfaces is even as to its highest divisions incomplete; and it is by no means easy to make it at once exhaustive and precise; an enumeration of all the prima facie possible cases would include forms which do not really exist. Thus the singular curve (if any) is of the order 1, 2, or 3; but in the case where the order is $= 3$, the curve, as is at once evident, cannot be a plane cubic, nor (among other excluded forms) a system of three non-intersecting lines. And certain forms of the singular curve, e.g.\ all but one of the admissible forms of a curve of the order 3, make the surface to be a scroll, so that, if (as is convenient) we wish to separate the scrolls, certain forms otherwise admissible must be excluded. The expression ``singular'' means double or cuspidal, or refers to a higher singularity, but the cases of higher singularity are very special. I will, at the cost of some inaccuracy, use the expression ``nodal'' as meaning, in general, double, but as including the signification ``cuspidal''; and, if there are any cases of higher singularity, as extending to cases of higher singularity: and I provisionally arrange the non-scrolar quartic surfaces as follows:
\begin{enumerate}[label=\arabic*.]
\item Without a nodal curve.
\item With a nodal line.
\item With a nodal conic, or line-pair (pair of intersecting lines).
\item With three nodal lines (not in the same plane) meeting in a point.
\end{enumerate}

The references, by the name of the author, and number (if any) of his paper, are to the subjoined list of Memoirs.

As to the scrolls, we have Cayley (3) and (4), and Cremona; the division into 12 species is, I believe, complete: see post, the remarks upon Schwarz's paper on quintic scrolls.

As regards the non-scrolar surfaces:

\textbf{1.}\quad Without a singular curve. The surface may be without a cnicnode (conical point), or it may have any number of cnicnodes up to 16, Cayley (7): the cases of singularity higher than a cnicnode are probably very numerous, but they have been scarcely at all examined. The memoir just referred to relates chiefly to the several cases of not more than 10 nodes; the cases of 11, 12, 13, 14, 15, 16 nodes are considered incidentally, Kummer (2), but it was not the object of his paper to make an enumeration, and there may be cases which are not considered; the discussion of the cases considered is very full and interesting. The case of 16 nodes is also considered, Kummer (1). As to the surface with 16 nodes, it is to be remarked that the wave-surface, or generally the surface obtained by the homographic deformation of the wave-surface called, Cayley (1), the ``tetrahedroid'' is a special form of surface with 16 nodes: its relation to the general surface is explained, Cayley (2).

\textbf{2.}\quad Quartic surface with nodal line: considered incidentally, Clebsch (2) and (3). There are through the nodal line 4 (in general different) quadric cones, having each of them a contact of the second order with the surface along the nodal line; any plane through the nodal line meets the surface in a conic having its two intersections with the nodal line for a pair of conjugate points in regard to the section by the plane of any one of the four quadric cones. And in particular, the tangent plane at any point of the nodal line is a plane through the nodal line which is a double tangent plane of the surface; viz., the tangent plane at any point of the nodal line meets the surface in the nodal line twice and in a conic touching the nodal line at the point. The cases of 1, 2, 3, 4 cnicnodes on the surface are considered, Korndorfer.

\textbf{3.}\quad Quartic surfaces with nodal conic. Such a surface may be without cnicnodes, or it may have 1, 2, 3, or 4 cnicnodes. The theory depends on that of the cubic surface; viz., taking any one of the 16 lines of the cubic surface, the scroll generated by the line in its motion along the line is a quartic surface having the line for nodal conic; that is, the locus of a line which meets each of two given non-intersecting lines, and also a given cubic curve, is a quartic surface having the line-pair (the two given lines) for a nodal conic and to each other. And the ground hereof appears, Geiser; viz., it is shown that the quartic surface with the nodal conic, is rationally transformable into a cubic surface, the 16 lines and the nodal conic corresponding respectively to the 16 lines and the conic of the cubic surface.

The several cases of 1, 2, 3, and 4 cnicnodes are considered, Korndorfer.

In the case where the nodal conic is the circle at infinity, the surfaces have been termed ``anallagmatic'' (perhaps ``bicircular'' would be a more convenient name), and a great deal has been written upon these surfaces by Moutard, Clifford, and others. Such a surface may of course have 1, 2, 3, or 4 cnicnodes; these surfaces, viz.\ the cnicnodal anallagmatics, in fact arise from the inversion of a quadric surface by the method of reciprocal radius vectors (that is, by the change of $x$, $y$, $z$ into $\frac{x}{r^2}$, $\frac{y}{r^2}$, $\frac{z}{r^2}$); the centre of inversion is a node on the quartic surface. If the quadric surface is a cone, there is another node, the inverse point of the vertex; if the quadric surface is one of revolution, there are two other nodes; and if it is a cone of revolution, there are three other nodes viz., in all, four nodes. The last-mentioned surface is, or includes, the Cyclide; viz., this is a quartic surface having the circle at infinity for a nodal curve, and having besides four nodes, which are a system of skew antipoints. The surface was first considered by Dupin (\textit{Applications de G\'eom\'etrie}, 1822); it has been discussed and its theory completed by Maxwell (Quart.\ Math.\ Journ.\ t.~ix.\ pp.~111--126, 1868) who has also considered the particular cases; and it was considered under a more general point of view by Moutard, and more recently by Cremona (in a Memoir ``Sulle transformazioni razionali nello spazio,'' Annali di Matematica, t.~V.\ 1868) who showed that the cyclide is the inverse of a quadric surface, and is also the envelope of a sphere which cuts at right angles each of four given spheres; or, what is the same thing, that the surface is the locus of a point whose distances from the centres of the four given spheres satisfy a linear equation; that is, the centres of the four given spheres are the foci of the surface. This theory includes of course the particular case where the four spheres reduce themselves to points, viz., the cyclide becomes the ``sphero-quartic'' or cyclide with four concyclic foci, the locus of a point the sum of whose distances from two given points is in a constant ratio to the sum of its distances from two other given points; or, more generally, the locus of a point such that $l\rho + m\rho' + n\rho'' = 0$; viz., this is the ``Cartesian,'' which may therefore be regarded as a particular case of the cyclide. In reference to the plane sections of the conic torus and its various particular cases, see De la Gournerie (2); the ordinary torus has been the subject of numerous papers by Darboux and others, and possesses very interesting properties.

In connexion with the foregoing, I speak of the surfaces having a cuspidal conic: the general case is briefly referred to, Cayley (6); viz., this is the surface (AA) the equation of which is $V^2 - a^2 u^2 = 0$, and which it is shown has a reciprocal of the order 6. A special case is the surface (AB) having a reciprocal of the order 3; viz., the quartic surface is here the reciprocal of the cubic surface XX = $12 - f^2 g$, Cayley (5). And it appears from the memoir last referred to, that there is another cubic surface, XVII = $12 - 25\beta - C_2$, the reciprocal of which is a quartic surface having a cuspidal conic. But the theory of the quartic surfaces with a cuspidal conic has been hardly at all considered.

I do not know that anything has been done in regard to the quartic surfaces where the nodal conic be a line-pair (pair of intersecting lines), a form which is admissible, but which probably belongs to the theory of surfaces with a nodal line.

\textbf{4.}\quad Quartic surface with three nodal lines (not in the same plane) meeting in a point. I have not met with any mention of such surfaces.

\vspace{0.5em}
\textit{Quartic Scrolls.}

As regards the scrolls, we have, as already mentioned, the division into 12 species, each of them a rational scroll, viz., the sections by a plane being in every case unicursal curves. The division was made according to the nature of the singularities (nodal curve, and cnicnodes, if any); but Schwarz's memoir effects the division according to the value of the ``deficiency'' of the surface, viz., the plane section is a quartic curve, and the deficiency of the surface is the deficiency of this quartic curve. Now the nodal curve of the quartic scroll is a line, and the plane section has accordingly a node; the plane section is thus a quartic curve with a node, and the deficiency is thus $= 1$ or $= 0$, according as the quartic curve is or is not a unicursal curve. But the plane section of a scroll is in every case a unicursal curve; and the deficiency is thus $= 0$ for a quartic scroll with a nodal line; for a quartic scroll without a nodal line the deficiency is $= 1$. But the case of a quartic scroll without a nodal line is a scroll which has the line at infinity for a nodal line, that is, a cylindric quartic scroll. The foregoing relates of course to the scrolls above referred to; viz., for a quartic scroll the deficiency is either 0 or 1; and of the 12 species, there are 10 for which the deficiency is $= 0$ (or which are unicursal), and 2 for which the deficiency is $= 1$. And this is the principle of classification in Schwarz's memoir; viz., for a quintic scroll the deficiency is $= 0$, 1, or 2; the number of species established being 10, 4, and 1 for these deficiencies respectively.

\vspace{0.5em}
\textit{List of Memoirs.}

\textbf{Cayley.}\quad
1. Sur la surface des ondes. Liouv.\ t.~xi.\ (1846), pp.~291--296, [47].

2. Sur un cas particulier de la surface du quatrieme ordre avec seize points singuliers. Crelle, t.~LXV.\ (1866), pp.~284--290, [356].

3. Second Memoir on Skew Surfaces, otherwise Scrolls. Phil.\ Trans., vol.~CLIV.\ (1864), pp.~559--577, [340].

4. Third Memoir on Skew Surfaces, otherwise Scrolls. Phil.\ Trans., vol.~CLIX.\ (1869), pp.~111--126, [410].

5. Memoir on Cubic Surfaces. Phil.\ Trans., vol.~CLIX.\ (1869), pp.~231--326, [412].

6. On the Quartic Surfaces $(U, V, W)^2 = 0$. Quart.\ Math.\ Journ.\ vol.~x.\ (1868), pp.~24--34; and vol.~xi.\ (1870), pp.~15--25, and pp.~111--113.

7. A Memoir on Quartic Surfaces. Proc.\ Lond.\ Math.\ Soc., vol.~III.\ (1870), pp.~19--69, [445].

\textbf{Clebsch.}\quad
1. Ueber die Fl"achen vierter Ordnung welche eine Doppelcurve zweiten Grades besitzen. Crelle, t.~LXIX.\ (1868), pp.~355--358.

2. Intorno alla rappresentazione di superficie algebriche sopra un piano. Atti del R.~Ist.~Lomb.\ (12 Nov.~1868), 13~p.

3. Ueber die Abbildung algebraischer Fl"achen, insbesondere der vierten und f"unften Ordnung. G"ottingen (1868), 42~p.;
``dritter Ordnung,'' Crelle, t.~LXV.\ (1866), pp.~359--380.

4. Ueber die ebene Abbildung der geradlinigen Fl"achen vierter Ordnung. G"ottingen (1868), 14~p.

5. Ueber die Abbildung einer Classe von Fl"achen f"unfter Ordnung. G"ott.~Abh., t.~XV. (1868).

\textbf{Cremona.}\quad
Sopra alcune superficie rationali. Atti (1868).

\textbf{De la Gournerie.}\quad
1. Note sur la cyclide. Comptes Rendus, t.~LXVII. (1868), pp.~1021--1024.

2. 'Etudes sur la cyclide. Journ.~de l''Ecole Polytechnique, t.~XXXV. (1869), pp.~81--94.

\textbf{Geiser.}\quad
Ueber die Fl"achen vierten Grades welche eine Doppelcurve zweiten Grades haben. Crelle, t.~LXX.\ (1869), pp.~249--257.

\textbf{Korndorfer.}\quad
[There was only a general reference; the several memoirs are:

1. Die Abbildung einer Fl"ache vierter Ordnung mit einer Doppelcurve zweiten Grades und einem oder mehreren Knotenpunkten. Math.~Ann., t.~I.\ (1869), pp.~592--626.

2. Fortsetzung dieses Aufsatzes. t.~II.\ (1870), pp.~41--64.

3. Ueber diejenigen Raumcurven deren Coordinaten sich als rationale Functionen eines Parameters darstellen. t.~III.\ (1871), pp.~415--423.

4. Die Abbildung einer Fl"ache vierter Ordnung mit zwei sich schneidenden Doppelgeraden. t.~III.\ (1871), pp.~496--522.

5. Die Abbildung einer Fl"ache vierter Ordnung mit einer Doppelcurve zweiten Grades welche aus zwei sich schneidenden unendlich nahen Geraden besteht. t.~IV.\ (1871), pp.~117--134.]

\textbf{Kummer.}\quad
1. {Surfaces of the fourth order with sixteen conical points.} Berl.~Monatsb.\ (1864), pp.~246--260, and 495--499.

2. Ueber die algebraischen Strahlensysteme, insbesondere "uber die der ersten und zweiten Ordnung. Berl.~Abh.\ (1866), pp.~1--120.

3. Ueber die Fl"achen vierten Grades auf welchen Schaaren von Kegelschnitten liegen. Berl.~Monatsb.\ (July, 1868). Crelle, t.~LXIV.\ (1864), pp.~66--76.

\textbf{Maxwell.}\quad
On the Cyclide. Quart.\ Math.\ Journ., t.~ix.\ (1867), pp.~111--126.

\textbf{Schwarz.}\quad
Ueber die geradlinigen Fl"achen f"unften Grades. Crelle, t.~LXVII.\ (1867), pp.~23--57.

\clearpage

%% ========================================================================
%% PAPER 450: Note on the Theory of the Rational Transformation Between Two Planes
%% ========================================================================
\section*{\textbf{450.}}
\addcontentsline{toc}{section}{450. Note on the Theory of the Rational Transformation Between Two Planes, and on Special Systems of Points}

\begin{center}
\textbf{NOTE ON THE THEORY OF THE RATIONAL TRANSFORMATION BETWEEN TWO PLANES, AND ON SPECIAL SYSTEMS OF POINTS.}
\end{center}

[From the \textit{Proceedings of the London Mathematical Society}, vol.~iii (1869--1871), pp.~196--198. Read December 8, 1870.]

In Prof.~Cremona's theory of the transformation of plane curves, the fundamental equations are taken to be
\begin{align}
\alpha_1 + 4\alpha_2 + 9\alpha_3 + \ldots &= n^2 - 1 \tag{1},\\
\alpha_1 + 3\alpha_2 + 6\alpha_3 + \ldots &= \tfrac{1}{2}(n^2 + 3n) - 2 \tag{2};
\end{align}
and from these we have as a consequence
\[
\alpha_1 + 3\alpha_2 + \ldots = \tfrac{1}{2}(n-1)(n-2) \tag{3};
\]
viz., the first equation expresses that any two curves of the system intersect in a single variable point; the second, that the curves form a r\'eseau, or system containing two arbitrary parameters; and the third, that the curves are unicursal.

In the equivalent theory of the rational transformation between two planes, as given in my ``Memoir on the Rational Transformation between Two Spaces,'' [447], we have the equation (1); but instead of the equation (2), it would prim\^a facie appear to be sufficient if we had the inequation
\[
\alpha_1 + 3\alpha_2 + 6\alpha_3 + \ldots > \tfrac{1}{2}(n^2 + 3n) - 2;
\]
but on the ground there explained, the case
\[
\alpha_1 + 3\alpha_2 + 6\alpha_3 + \ldots < \tfrac{1}{2}(n^2 + 3n) - 2
\]
is excluded, and we thus have the equation (2), giving with (1) the equation (3).

I believe the better course is to assume (1) and (3) as the fundamental equations, from them deducing (2); and we thus also get over a difficulty presently referred to, but which did not occur to me when the memoir was written.

In fact, starting with the equations $x' : y' : z' = X : Y : Z$ (which are to give $x : y : z = X' : Y' : Z'$), we have in the first instance the equation (1). Moreover, establishing for $x', y', z'$ a linear equation $ax' + by' + cz' = 0$, we have corresponding hereto a curve $aX + bY + cZ = 0$, and the coordinates $x, y, z$ of a point on this curve are proportional to $X' : Y' : Z'$; that is, substituting for $z'$ the value $-(ax' + by')/c$, they are proportional to rational and integral (homogeneous) functions of $(x', y')$, that is, to rational and integral functions of the single parameter $x' : y'$; wherefore the curve $aX + bY + cZ = 0$ is unicursal; whence the equation (3). The like change may be made in the theory of the rational transformation between two spaces; and it is in this case a more important one.

The difficulty is as follows: It is not self-evident that we are at liberty to assume
\[
\alpha_1 + 3\alpha_2 + 6\alpha_3 + \ldots > \tfrac{1}{2}(n^2 + 3n) - 2
\]
for imagine that we had a system of $(\alpha_1, \alpha_2, \alpha_3, \ldots)$ points, such that $\alpha_1 + 4\alpha_2 + \ldots$ being $= n^2 - 1$, and $\alpha_1 + 3\alpha_2 + \ldots$ being $> \tfrac{1}{2}(n^2 + 3n) - 2$, the points were such that the conditions in question (viz., the condition that the curve passes once through each of the points $\alpha_1$, twice through each of the points $\alpha_2$, \&c.) were not independent, but equivalent to a less number of conditions. Suppose, for instance, a quintic curve, $n = 5$; and consider the system of 12 points and 2 double points (12 points through which the curve passes once, and 2 points through which it passes twice); we have here
\[
\alpha_1 + 4\alpha_2 = 12 + 8 = 20 = n^2 - 1,\quad \alpha_1 + 3\alpha_2 = 12 + 6 = 18 = \tfrac{1}{2}(n^2 + 3n) - 2;
\]
so that the equations (1) and (2) are satisfied. But there is here no reduction; the 18 conditions are independent. But we may, it would appear, have 12 points and 2 double points such that, for a quintic curve passing once through each of the 12 points and twice through each of the 2 points, the number of conditions actually imposed (instead of being $12 + 3 \cdot 2 = 18$) is $= 17$.'' The construction is as follows: viz., starting with the 2 points and any 10 points, we may draw a quartic passing twice through the first of the 2 points, once through the second of them, and through the 10 points; and another quartic passing twice through the second of the 2 points, once through the first of them, and through the 10 points: the two quartics intersect in the 2 points each twice, in the 10 points, and in 2 new points, forming, with the 10 points, a system of 12 points; and the first-mentioned 2 points and the 12 points form the system in question.

A more complicated case, $\alpha_1 = 10$, $\alpha_2 = 6$, $\alpha_3 = 1$, occurs in Dr Noether's paper, ``Ueber Fl\"achen, welche Schaaren rationaler Curven besitzen,'' [Math.\ Ann., t.~iii.\ (1871), pp.~161--227]. Except these two, I do not know any other case of a special system for which $\alpha_2, \alpha_3, \ldots$ are not all $= 0$; the investigation of such systems would, I think, be very interesting.

[A concluding paragraph of seven lines gave some corrections to the ``Memoir on the Rational Transformation between Two Spaces,'' 447, which corrections are made in the present reprint of that paper.]

\clearpage

%% ========================================================================
%% PAPER 451: A Second Memoir on Quartic Surfaces
%% ========================================================================
\section*{\textbf{451.}}
\addcontentsline{toc}{section}{451. A Second Memoir on Quartic Surfaces}

\begin{center}
\textbf{A SECOND MEMOIR ON QUARTIC SURFACES.}
\end{center}

[From the \textit{Proceedings of the London Mathematical Society}, vol.~iii (1869--1871), pp.~198--202. Read December 8, 1870.]

In my Memoir on Quartic Surfaces, ante pp.~19--69, [445], although remarking (see No.~79) that the identification was not completely made out, I tacitly assumed that the symmetroid and the decadianome (each of them a quartic surface with 10 nodes) were in fact identical. There is yet a good deal which I cannot completely explain; but the truth appears to be, that the decadianome includes two cases of coordinate generality, say the sextic decadianome, and the bicubic decadianome $=$ symmetroid: viz., in the first of these the circumscribed cone, having for vertex any one of the 10 nodes, is a proper sextic cone with 9 double lines; in the second it is a system of two cubic cones, intersecting, of course, in 9 lines, which are double lines of the aggregate sextic cone: or, in the notation of the Table No.~11, in the case of the sextic decadianome, the circumscribed cones are each of them $6_9$; in that of the bicubic decadianome $=$ symmetroid, they are each of them $(3, 3)$. We thus arrive at a very remarkable system of 10 points in space, viz., giving the name ``ennead'' to any 9 points in piano, which are the intersections of two cubic curves, or to any 9 lines through a point which are the intersections of two cubic cones; the 10 points in space are such that, taking any one whatever of them as vertex, and joining it with the remaining 9, the 9 lines form an ennead. I purpose in the present short Memoir to consider the theories in question; the paragraphs are numbered consecutively with those of the Memoir on Quartic Surfaces.

\subsection*{Plane Sextic Curve with 9 Nodes.}

\textbf{\S 110.}\quad A sextic curve contains 27 constants; and the number of conditions to be satisfied in order that a given point may be a node is $= 3$. Hence it would at first sight appear that the curve could be found so as to have 9 given nodes; this would be $9 \times 3 = 27$ conditions, or the curve would be completely determinate. But observe that through the 9 given points we have a determinate cubic curve $U = 0$; we have therefore $U^2 = \text{a sextic curve}$, and the only sextic curve with the 9 given nodes; that is, there is not in a proper sense any sextic curve with the 9 given nodes. The number of given nodes is thus $= 8$ at most.

\textbf{\S 111.}\quad The sextic curve with 8 given nodes should contain $27 - 3 \cdot 8 = 3$ constants. We may through the 8 given points draw the two cubics $P = 0$, $Q = 0$; and we have then $(a, b, c \mid P, Q)^2 = 0$, a bicubic, or improper sextic curve having the 8 nodes, and also a ninth node, viz., the remaining point of intersection of the two cubic curves, or say the remaining point of the ennead. Hence if $V = 0$ be any particular sextic curve having the 8 given nodes, we have
\[
(a, b, c \mid P, Q)^2 + \theta V = 0
\]
a proper sextic curve having the 8 given nodes; and this, as containing the right number ($= 3$) of constants, will be the general sextic curve having the 8 given nodes.

\textbf{\S 112.}\quad There will be a ninth node if $\theta = 0$; viz., the curve is then $(a, b, c \mid P, Q)^2 = 0$, a bicubic, or improper sextic curve, having for nodes the 9 points of the ennead. Observe that the ninth node is here a point completely and uniquely determined by means of the given 8 nodes. Moreover the number of constants is $= 2$, so that we have here a general (improper) solution of the question of finding a sextic curve with 9 nodes, 8 of them given.

\textbf{\S 113.}\quad But if $\theta$ is not $= 0$, then the ninth node must be a point on the curve $J(P, Q, V) = 0$; viz., this is a curve of the order 9, determined by means of the given 8 points; say it is the ``dianodal curve'' of these 8 points, and, as is easy to see, it has each of these 8 points for a node. The ninth node of the sextic may be any point whatever on the dianodal curve; and regarding it as a given point, the sextic will still contain 1 constant; that is, we have the general solution of the problem of finding a sextic curve with 9 nodes, 8 of them given, and the 9th a given point on the dianodal curve.

\textbf{\S 114.}\quad So long as the 8 points are arbitrary, the dianodal curve does not pass through the 9th point of the ennead, and the two cases above considered are mutually exclusive. It will be noticed how closely analogous this theory of the plane sextic with 9 nodes, is to that of the quartic surface with 8 nodes.

\textbf{\S 115.}\quad Of course, instead of the plane sextic curve, we may have a sextic cone; such a cone has at most 8 given double lines; and if there be a 9th double line, then there are the two cases of coordinate generality; viz., (1), the new double line is the ninth line of the ennead, the cone being in this case not a proper sextic cone, but a bicubic cone; (2), the new double line may be any line through the vertex which lies on the dianodal cone of the 8 given double lines.

\subsection*{Each circumscribed cone of the Symmetroid is $(3, 3)$.}

\textbf{\S 116.}\quad Using $(x, y, z, w)$ as current coordinates of a point of the symmetroid, I take $S$, $T$, $U$, $V$ to be quadric functions of the coordinates $(\alpha, \beta, \gamma, \delta)$; the equation of the symmetroid is therefore given by
\[
xS + yT + zU + wV = \text{cone},
\]
and the nodes thereof are determined by
\[
xS + yT + zU + wV = \text{plane-pair}.
\]
Suppose that a node is $(x = 0, y = 0, w = 0)$; the condition for this is $V = \text{plane-pair}$; and we may without loss of generality write $V = \gamma^2 + \delta^2$. Hence, putting for shortness $\Omega = xS + yT + zU$, that is a quadric function $(a, b, c, d, f, g, h, l, m, n \mid \alpha, \beta, \gamma, \delta)^2$, wherein the coefficients $a, b, \ldots$ are arbitrary linear functions of $(x, y, z)$, but not containing $w$, the equation of the symmetroid is given by
\[
\Omega + w(\gamma^2 + \delta^2) = \text{cone}.
\]

\textbf{\S 117.}\quad It follows that the equation is
\[
\begin{vmatrix}
a, & h, & g \\
h, & b, & f \\
g, & f, & c + w
\end{vmatrix}
= 0;
\]
viz., this is
\[
V^2 + w(\partial_c V + \partial_d V)V + w^2 \partial_c V \cdot \partial_d V = 0,
\]
where $V$ denotes the foregoing determinant, writing therein $w = 0$. Or, observing that $V$ contains $c$ and $d$ each only linearly, the equation may be written
\[
V^2 + w(\partial_c + \partial_d)V + w^2 \partial_c \partial_d V = 0,
\]
which is a quartic surface having, as it should have, the point $(0, 0, 0)$ for one of its ten nodes.

\textbf{\S 118.}\quad The equation of the circumscribed cone is
\[
(\partial_c V - \partial_d V)^2 + 4(\partial_c V \cdot \partial_d V - V \cdot \partial_c \partial_d V) = 0.
\]
But we have identically
\[
\partial_c V \cdot \partial_d V - (\partial_c \partial_d V)^2 = V \cdot \partial_c \partial_d V;
\]
so that the equation is
\[
(\partial_c V - \partial_d V)^2 + (\partial_{cd} V)^2 = 0;
\]
a sextic cone breaking up into the two cubic cones
\[
\partial_c V - \partial_d V + i\partial_{cd} V = 0,\quad \partial_c V - \partial_d V - i\partial_{cd} V = 0,
\]
so that the cone is $(3, 3)$. And since clearly the point $(0, 0, 0)$ may be regarded as representing any one whatever of the 10 nodes, it follows that for any node whatever of the symmetroid, the circumscribed cone is $(3, 3)$, so that, as stated above, bicubic decadianome $=$ symmetroid.

\subsection*{Deductions from the foregoing theory.}

\textbf{\S 119.}\quad Referring to No.~85 of the original memoir, it appears that, with 6 given points as nodes, we can actually find for the symmetroid an equation containing 6 constants. I cannot discover any ground for doubting that 3 of these may be determined so as to give to the symmetroid a seventh given node; and I therefore assume that with 7 given points as nodes, an equation can be found with 3 constants. The symmetroid is certainly not octadic, hence the eighth node must lie on the dianodal surface of the 7 given points. I can discover no ground for doubting but that two of the constants may be determined so that the eighth node shall be any given point whatever on the dianodal surface of the 7 points; and (this being so) that further the remaining constant may be determined so that the ninth node shall be any given point whatever on the dianodal curve of the 8 points. But if all this be so, the consequence is very remarkable; the tenth node is not any one whatever of the 22 dianodal centres of the 9 points, but it is a uniquely determinate ``enneadic centre,'' viz., we must have the following theorem:

\textbf{\S 120.}\quad ``Take any 7 points; an eighth point at pleasure on the dianodal surface of the 7 points; a ninth point at pleasure on the dianodal curve of the 8 points. In the system of 9 points so determined, take any one as vertex, and joining it with the remaining 8, construct the ninth line of the ennead. Performing this construction with each of the 9 points successively as vertex, we obtain 9 lines passing through the 9 points respectively. These 9 lines meet in a point which is the `enneadic centre' of the 9 points: and further, the 10 points form a completely symmetrical system, so that each one of them is the enneadic centre of the remaining 9.''

\textbf{\S 121.}\quad Assuming that the 9 lines do intersect so as to give rise to an enneadic centre, there is no difficulty in conceiving that the loci, which by their intersection determine the dianodal centres, do each of them pass through the enneadic centre; so that this enneadic centre counts once or more among the dianodal centres, and the number of proper dianodal centres, instead of being $= 22$, will be suppose $= 22 - \omega$; and if, further, the 9 points, together with the enneadic centre, are the nodes of a symmetroid, but the 9 points together with any one of the $22 - \omega$ dianodal centres are the nodes of a sextic decadianome, then we must also have as follows:

\textbf{\S 122.}\quad ``Considering any 9 points as above; taking any one as vertex, and joining it with the remaining 8, these 8 lines determine a dianodal ninthic cone. We have thus 9 dianodal cones, which cones pass all of them through the same $22 - \omega$ points.''

\textbf{\S 123.}\quad I am not able to verify these theorems \textit{a posteriori}. It appears to me that the theorem in regard to the enneadic centre subsists for a system of 9 points such as referred to in the statement; but that if by possibility the statement be too general, the theorem must, at all events, subsist for a more special system of 9 points; and that there certainly exist systems of 10 points, such that each 9 of the points have as an enneadic centre the tenth point. (I have since ascertained that if a quartic surface with 10 nodes has a single node $(3, 3)$, the surface is a symmetroid; whence, by what precedes, the remaining nine nodes are each of them $(3, 3)$. Added 25 March, 1871.)

\textbf{\S 124.}\quad I notice, as a subject of investigation, the following system of correspondence viz., given any 8 points in space: then to every point in space corresponds a line through this point, viz., the ninth line of the ennead obtained by joining the point with the 8 given points respectively; and to each line in space a point or points on the line, viz., the point or points for each of which the line is the ninth line of the ennead obtained by joining the point with the eight given points respectively.

\clearpage

%% ========================================================================
%% PAPER 452: On an Analytical Theorem from a New Point of View
%% ========================================================================
\section*{\textbf{452.}}
\addcontentsline{toc}{section}{452. On an Analytical Theorem from a New Point of View}

\begin{center}
\textbf{ON AN ANALYTICAL THEOREM FROM A NEW POINT OF VIEW.}
\end{center}

[From the \textit{Proceedings of the London Mathematical Society}, vol.~iii (1869--1871), pp.~220, 221. Read February 9, 1871.]

The theorem is a well-known one, derived from the equation
\[
(az^2 + 2bz + c)w^2 + 2(a'z + 2b'z + c')w + a''z^2 + 2b''z + c'' = 0;
\]
viz., considering this equation as establishing a relation between the variables $z$ and $w$, and writing it in the forms
\[
\Omega = Aw^2 + 2Bw + C = A'z^2 + 2B'z + C' = 0,
\]
(where, of course, $A$, $B$, $C$ are quadric functions of $z$, and $A'$, $B'$, $C'$ quadric functions of $w$,) we have
\[
0 = \frac{dw}{\sqrt{B'^2 - A'C'}} + \frac{dz}{\sqrt{B^2 - AC}},\quad
\Omega = (Aw + B)\,dw + (A'z + B')\,dz;
\]
but in virtue of the equation $\Omega = 0$, we have $Aw + B = \sqrt{B^2 - AC}$, and $A'z + B' = \sqrt{B'^2 - A'C'}$, and the differential equation thus becomes
\[
\frac{dw}{\sqrt{B'^2 - A'C'}} = \frac{dz}{\sqrt{B^2 - AC}},
\]
where $B'^2 - A'C'$ and $B^2 - AC$ are quartic functions of $w$ and $z$ respectively. This is, of course, integrable (viz., the integral is the original equation $\Omega = 0$); and it follows, from the theory of elliptic functions, that the two quartic functions must be linearly transformable into each other; viz., they must have the same absolute invariant $\frac{I^3}{J^2}$. It is, in fact, easy to verify, not only that this is so, but that the two functions have the same quadrinvariant $I$, and the same cubinvariant $J$.

The new point of view is, that we take the coefficients $a$, $b$, \&c., to be homogeneous functions of $(x, y)$, their degrees being such that the equation $\Omega = 0$ is a quartic equation $(\mid\mid\mid x, y, z, w)^4 = 0$: viz., this equation now represents a quartic surface having a node (conical point) at the point $(x = 0,\; y = 0,\; z = 0)$, and also a node at the point $(x = 0,\; y = 0,\; w = 0)$, say, these points are $O$, $O'$ respectively. The equation $B'^2 - A'C' = 0$ gives the circumscribed sextic cone having $O$ for its vertex, and the equation $B^2 - AC = 0$ the circumscribed sextic cone having $O'$ for its vertex; each of these cones has the line $OO'$ ($x = 0,\; y = 0$) for a nodal line, as appears geometrically, and also by the equations containing $z$, $w$ respectively in the degree 4. Considering $B'^2 - A'C'$ as a quartic function of $z$, its quadrinvariant is a function $(x, y)^4$, and its cubinvariant a function $(x, y)^6$; and similarly, considering $B^2 - AC$ as a quartic function of $w$, its invariants are functions $(x, y)^4$ and $(x, y)^6$. We have thus, between the two cones, a geometrical relation answering to the analytical one of the identity of the invariants; but the nature of this geometrical relation is not obvious; and it presents itself as an interesting subject of investigation.

\clearpage

%% ========================================================================
%% PAPER 453: On a Problem in the Calculus of Variations
%% ========================================================================
\section*{\textbf{453.}}
\addcontentsline{toc}{section}{453. On a Problem in the Calculus of Variations}

\begin{center}
\textbf{ON A PROBLEM IN THE CALCULUS OF VARIATIONS.}
\end{center}

[From the \textit{Proceedings of the London Mathematical Society}, vol.~iii (1869--1871), pp.~221, 222. Read February 9, 1871.]

The problem is, $z = \frac{1}{3}(3x - y^2)y$, to find $y$ a function of $x$ such that $\int z\,dx = \max.$ or $\min.$, subject to a given condition $\int y\,dx = c$ (the limits of each integral being $x_0$, $x_1$, where these quantities are each positive, and $x_1 > x_0$). The ordinary method of solution gives $y^2 = x + \lambda$, where $(x_1 + \lambda)^2 - (x_0 + \lambda)^2 = \frac{3}{2}c$; so long as $c$ is not less than $(x_1 - x_0)$, there is a real value of $\lambda$, but for a smaller value of $c$ there is no real value. The difficulty arising in this last case is somewhat illustrated by replacing the original problem by a like problem of ordinary maxima and minima; viz., $x_1$, $x_2$, $\ldots$, $x_n$ being given positive values of $x$, in the order of increasing magnitude; and if, in general, $z_i = \frac{1}{3}(3x_i - y_i^2)y_i$, then the problem is to find $y_i$ a function of $x_i$, such that $\Sigma z_i = \max.$ or $\min.$, subject to the condition $\Sigma y_i = c$. We have here $y_i^2 = x_i + \lambda$ where $\lambda$ is to be determined by the condition $\Sigma y_i = c$; the remainder of the investigation turns on the question of the sign $y_i = +\sqrt{x_i + \lambda}$ or $y_i = -\sqrt{x_i + \lambda}$, to be taken for the several values of $i$ respectively.

\clearpage

%% ========================================================================
%% PAPER 454: A Third Memoir on Quartic Surfaces
%% ========================================================================
\section*{\textbf{454.}}
\addcontentsline{toc}{section}{454. A Third Memoir on Quartic Surfaces}

\begin{center}
\textbf{A THIRD MEMOIR ON QUARTIC SURFACES.}
\end{center}

[From the \textit{Proceedings of the London Mathematical Society}, vol.~iii (1869--1871), pp.~234--266. Read April 13, 1871.]

The present Memoir is a continuation of my former researches on Nodal Quartic Surfaces, [445, 451]. The leading idea is, that for a quartic surface with $k$-nodes, given the nature of the circumscribed, $(k - 1)$ nodal, sextic cone belonging to any one node of the surface (for instance, $k = 10$, that it is a cone $(3, 3)$ composed of two cubic cones), we thereby determine the equation of the quartic surface, and consequently the nature of the remaining $(k - 1)$ nodes thereof. By means of this general theory I complete, in an essential point, the theory of the Symmetroid; viz., I show that a 10-nodal quartic surface having a single node $(3, 3)$ is a Symmetroid; whence, as appears by my second Memoir, [451], each of the remaining nine nodes is also a node $(3, 3)$; and we have the theory of the remarkable system of ten points in space such that, joining any one of them with the remaining nine, the nine lines thus obtained are the intersections of two cubic cones. A large part of the Memoir is devoted to the consideration of the surfaces with 16, 15, 14, and 13 nodes: this is substantially a reproduction of the results obtained by Kummer in the Memoir ``Ueber die algebraischen Strahlensysteme, \&c.,'' already referred to; but the results in question are brought into connexion with the theory of the present Memoir, and they are, by a change of the constants, exhibited in a form of much greater symmetry and elegance. I attach importance also to the square diagrams by means of which I have exhibited, in a compendious form, the relation between the several nodes and circumscribed sextic cones.

The paragraphs are numbered consecutively with those of the first and second Memoirs.

\subsection*{Preliminary Considerations and Classification.}

\textbf{\S 125.}\quad I call to mind that if a quartic surface has a node (conical point), then there is for this node a tangent quadricone and a circumscribed sextic cone; viz., if the surface has $(k - 1)$ other nodes, or in all $k$ nodes, then the sextic cone has $(k - 1)$ nodal lines (passing through the other nodes respectively), and we have thus for the different forms of the sextic the table No.~11; viz., this is

\begin{center}
\begin{tabular}{|c|c|c|}
\hline
\multicolumn{2}{|c|}{Nodes of} & Circumscribed Sextic Cone. \\
\cline{1-2}
Surface. & Nodal Lines of Cones. & \\
\hline
1 & 0 & $6$ \\
2 & 1 & $6_1$ \\
3 & 2 & $6_2$ \\
4 & 3 & $6_3$ \\
5 & 4 & $6_4$ \\
6 & 5 & $6_5$ \\
7 & 6 & $6_6$ \\
8 & 7 & $6_7$ \\
9 & 8 & $6_8$, $5_1$, $4_2$, $2$ \\
10 & 9 & $6_9$, $5_1$, $4_2$, $2$, $4_1$, $1$, $1$, $3_1$, $3$ \\
11 & 10 & $6_{10}$, $5_1$, $4_2$, $2$, $4_1$, $1$, $1$, $3_1$, $3_1$, $3$, $2$, $1$ \\
12 & 11 & $5_1$, $4_2$, $2$, $4_1$, $1$, $1$, $3_1$, $3_1$, $3$, $2$, $1$ \\
13 & 12 & $4_1$, $1$, $1$, $3_1$, $2$, $1$, $3$, $1$, $1$, $1$, $2$, $2$, $2$ \\
14 & 13 & $3_1$, $1$, $1$, $1$, $2$, $2$, $1$, $1$ \\
15 & 14 & $2_1$, $1$, $1$, $1$, $1$ \\
16 & 15 & $1$, $1$, $1$, $1$, $1$, $1$ \\
\hline
\end{tabular}
\end{center}

viz., $6$ denotes a proper sextic cone without nodal lines; $6_1$ a proper sextic cone with one nodal line; $5$, $1$ a proper quintic cone and a plane, \&c.

We may distinguish the nodes according to the sextic cones; thus, a node $6$ means a node for which the circumscribed cone is a proper sextic cone, $(1, 1, 1, 1, 1, 1)$ a node where the circumscribed cone breaks up into six planes, \&c.

\textbf{\S 126.}\quad A 16-nodal surface has 16 nodes $(1, 1, 1, 1, 1, 1)$, and a 15-nodal surface has 15 nodes $(2, 1, 1, 1, 1)$; but, for a 14-nodal surface, the question arises how many nodes are $(3_1, 1, 1, 1)$, and how many $(2, 2, 1, 1)$. It was remarked, No.~13, that the only possible cases were 14, 0; 8, 6; or 2, 12; and that we might, in like manner, limit the number of possible cases for other values of $k$; but that the inquiry was not then further pursued. I resume this inquiry, but without obtaining as yet a complete answer.

\textbf{\S 127.}\quad It is to be observed that a line joining any two nodes is not, in general, a line on the surface, but that it may be so; the surfaces for which this is so (viz., any surface which contains upon it a line through two nodes) form, however, a class by themselves, which at present I altogether exclude from consideration. This being so, it will appear in the sequel that there is but one kind of surface having a node $(2, 2, 1, 1)$, and but one kind of surface having a node $(3_1, 1, 1, 1)$. Now there is a surface, Kummer's 14-nodal, the nodes of which are $8(3_1, 1, 1, 1) + 6(2, 2, 1, 1)$; wherefore the two kinds are identical, and are each of them Kummer's 14-nodal surface. Similarly, for the 13-nodal surfaces, there is but one kind having a node $(4_1, 1, 1)$, but one kind having a node $(3, 1, 1, 1)$, and moreover but one kind having a node $(3_1, 2, 1)$; and we have Kummer's 13-nodal surface with the nodes $3(4_1, 1, 1) + 1(3, 1, 1, 1) + 9(3_1, 2, 1)$; hence the three kinds are identical with each other and with Kummer's. Moreover, there is but one kind having a node $(2, 2, 2)$; hence all the other nodes must be $(2, 2, 2)$, and we have a surface $13(2, 2, 2)$ not given by Kummer. And in like manner for the 12-nodal surfaces, we have the two kinds given by Kummer, and a third kind $12(4_1, 1, 1)$ not given by him; the arrangement thus far being

\begin{center}
\begin{tabular}{c|l}
No.\ of Nodes. & Character of Surface. \\
\hline
16 & $16(1, 1, 1, 1, 1, 1)$, \\
15 & $15(2, 1, 1, 1, 1)$, \\
14 & $8(3_1, 1, 1, 1) + 6(2, 2, 1, 1)$, \\
13 ($\alpha$) & $3(4_1, 1, 1) + 1(3, 1, 1, 1) + 9(3_1, 2, 1)$, \\
\quad ($\beta$) & $13(2, 2, 2)$, \\
12 & ($\alpha$) $12(4_1, 2)$, \\
\quad ($\beta$) & $2(5_1, 1) + 6(3_1, 3_1) + 4(3, 2, 1)$, \\
\quad ($\gamma$) & $12(4_1, 1, 1)$.
\end{tabular}
\end{center}

\textbf{\S 128.}\quad But in the next following case we have Kummer's surface, viz.
\[
11(\alpha)\quad 1(6_{10}) + 10(3_1, 3),
\]
and I do not know whether one, two, or three kinds of surface having nodes $(4_1, 1, 1)$, $(4_1, 2)$, and $(5_1, 1)$. And in the next case we have (as will appear) the Symmetroid, viz.,
\[
10(\alpha)\quad 10(3, 3),
\]
and I do not know how many kinds of surfaces having a node or nodes $6_{10}$, $(5_1, 1)$, $(4_1, 2)$, $(4_1, 1, 1)$.

It will be observed that the present division has nothing to do with the octadic and disinodal division in the former Memoir.

\textbf{\S 129.}\quad I consider a conic $A = 0$, and any six tangents thereof, $t_1 = 0$, $t_2 = 0$, $t_3 = 0$, $t_4 = 0$, $t_5 = 0$, $t_6 = 0$; we have an identical equation which might be written $AC - B^2 = t_1t_2t_3t_4t_5t_6$, but it will be convenient, introducing a constant factor $K$, to write it
\[
AC - B^2 = Kt_1t_2t_3t_4t_5t_6,
\]
$B$ being a cubic function and $C$ a quartic function of the coordinates.

Consider now the series of factors, such as
\begin{align*}
lA &+ mt_1t_2, \\
sA &+ mt_1t_2t_3, \\
UA &+ mt_1t_2t_3t_4, \\
VA &+ mt_1t_2t_3t_4t_5, \\
WA &+ mt_1t_2t_3t_4t_5t_6,
\end{align*}
where $m$ is a constant, $l$ a constant, $s$ a linear function, $U$, $V$, $W$ functions of the degrees 2, 3, 4 respectively; and compose with one or more such factors an expression involving the term $Kt_1t_2t_3t_4t_5t_6$; for instance, such an expression is
\[
(sA + mt_1t_2t_3)(l'A + m't_4t_5t_6);
\]
this is, of the form $An + Kt_1t_2t_3t_4t_5t_6$, viz., $An + (AC - B^2)$, or $A(n + C) - B^2$, say $A\Gamma - B^2$; or what is the same thing, introducing a new coordinate $w$, we have a quadric function
\[
Aw^2 + 2Bw + \Gamma,
\]
$AV - B^2$ --- the discriminant of which, $B^2 - A\Gamma$, is equal to the expression in question.

\textbf{\S 130.}\quad In the sequel $(x, y, z, w)$ are considered as the coordinates of a point in space; $A = 0$ is thus a quadric cone, $t_1 = 0$, $t_2 = 0$, $\ldots$, $t_6 = 0$, any six tangent planes thereof; and hence $Aw^2 + 2Bw + \Gamma = 0$ a quartic surface, having the point $(x = 0,\; y = 0,\; z = 0)$ for a node, whereof the circumscribed cone $A\Gamma - B^2 = 0$ breaks up in the assumed manner.

Thus, in order that the circumscribed cone may be as above
\[
(sA + mt_1t_2t_3)(l'A + m't_4t_5t_6),
\]
we have only to assume
\[
\Gamma = C + mm'(sl'A + sm't_4t_5t_6 + l'mt_1t_2t_3),
\]
and so in other cases. Observe that $sA + mt_1t_2t_3 = 0$ is a cubic cone, which, so long as $m$, $s$, are arbitrary, has no nodal line; but establishing a single relation (say $s$ remains arbitrary, but a proper value is assigned to $m$) it will be a cubic cone having a nodal line. And so $UA + mt_1t_2t_3t_4 = 0$ is a quartic cone without any nodal line, but by particularising the constants it may be made to have one, two, or three nodal lines. Such nodal determinations are obviously required in order that the formula may extend to all the before-mentioned forms of the circumscribed cone. The foregoing analysis is the foundation of the whole theory: I have given it, as above, apart from the theory, in order that the nature of it may be the better perceived; but I have now to bring it into connexion with the theory.

\subsection*{On the Sextic Curves, $AC - B^2 = 0$.}

\textbf{\S 131.}\quad I revert to the consideration of plane curves. The equation of a sextic curve $(\mid\mid\mid x, y, z)^6 = 0$ cannot be in general expressed in the form $AC - B^2 = 0$, where the degrees of $A$, $B$, $C$ are 2, 3, 4 respectively; in fact, the existence of such a form implies that there is a conic $A = 0$ touching the sextic 6 times; and since a conic can only be made to satisfy 5 conditions, there is not in general any such conic.

\textbf{\S 132.}\quad Such conic, when it exists, is said to be inscribed in the sextic, and the sextic to be circumscribed about the conic, or to be an ``amphigram;'' and then, $A = 0$ being the equation of the conic, that of the sextic is expressible in the form in question $AC - B^2 = 0$. It is clear that $B = 0$ is a cubic curve passing through the 6 points of contact of the conic with the sextic, and that any such curve may be taken for the cubic $B$; in fact, if $B = B'$ is any particular solution of the equation $AC - B^2 = 0$, but changing the cubic $B'$ by the addition of a linear function of the coordinates, the equation of the general cubic is $B = B' + pA = 0$; and then writing
\[
A = A,\quad B = B' + pA,\quad C = C + 2B'p + Ap^2,
\]
we have $AC - B^2 = AC - B'^2$; so that the original form $AC - B'^2 = 0$ becomes $AC - B^2 = 0$. But the cubic $B = 0$ being assumed at pleasure, the quartic $C = 0$ is a determinate curve.

\textbf{\S 133.}\quad It is to be observed that a sextic curve may be an amphigram in more than one way: certainly in two, three, or four, and possibly in a greater number of ways. For the equation of the curve contains 27 constants, and hence determining the sextic so as to touch 4 given conics each of them 6 times, there are still 3 constants and the curve will be an amphigram in regard to each of the 4 conics; say it is a quadruple amphigram. But in the sequel we are only concerned with a sextic curve considered as an amphigram in regard to a given conic $A = 0$ (no attention being paid to the other inscribed conics, if any); and then, by what precedes, taking $B = 0$ any cubic whatever through the 6 points of contact, we have a determinate quartic curve $C = 0$, and the equation of the sextic curve assumes the form $AC - B^2 = 0$.

\textbf{\S 134.}\quad The curves $A = 0$, $B = 0$, $C = 0$ contain respectively 5, 9, 14 constants; whence considering the function $B$ as containing an arbitrary constant factor, for the curve $AC - B^2 = 0$, the number of constants is prim\^a facie $5 + 9 + 14 + 1 = 29$; but on account of the arbitrary linear function $p$, the real number is $29 - 3 = 26$: this is right, for a sextic curve contains 27 constants; and the curve being an amphigram, there is one relation between the constants, $27 - 1 = 26$.

\textbf{\S 135.}\quad Suppose now that the sextic curve $AC - B^2 = 0$ breaks up into two or more separate curves, say into the two curves $P = 0$, $Q = 0$ of the orders $f$, $g$ respectively, $f + g = 6$. We have
\[
AC - B^2 = PQ = 0;
\]
and the conic $A = 0$ touching the sextic six times, must, it is clear, touch the curves $P = 0$, $Q = 0$, $f$ and $g$ times respectively. And so when the sextic breaks up into any number of curves, each component curve $P = 0$ of the order $f$ must touch the sextic $f$ times.

\textbf{\S 136.}\quad It follows that if the sextic break up into six lines, say $AC - B^2 = t_1t_2t_3t_4t_5t_6 = 0$, then that each of the lines $t_1 = 0$, $t_2 = 0$, $\ldots$, $t_6 = 0$ is a tangent to the conic. And conversely, starting with the conic $A = 0$ and any six tangents thereof $t_1 = 0$, $t_2 = 0$, $\ldots$, $t_6 = 0$, we have an identity of the form in question. In fact, taking any two of the tangents, say $t_1 = 0$ and $t_2 = 0$, then, if $p = 0$ be the equation of the line joining their points of intersection, the equation of the conic will be of the form $t_1t_2 + p^2 = 0$, that is, we may write $A = t_1t_2 + p^2$, or what is the same thing, $t_1t_2 = A - p^2$. (Considering $A$ as a given quadric function of the coordinates, this of course implies that the implicit constant factors of $t_1$, $t_2$, $p$ are properly determined.) Similarly, $q = 0$ being the line through the points of contact of $t_3$, $t_4$, and $r = 0$ that through the points of contact of $t_5$, $t_6$, we have $t_3t_4 = A - q^2$ and $t_5t_6 = A - r^2$; whence, to satisfy the equation $AC - B^2 = t_1t_2t_3t_4t_5t_6$, we have only to assume $B = lA + pqr$, $l$ an arbitrary linear function of the coordinates, and the equation then gives
\[
C = l^2A^2 - A(p^2 + q^2 + r^2) + (p^2q^2 + r^2p^2 + p^2r^2) + l^2 \cdot \ldots + 2lpqr.
\]

\textbf{\S 137.}\quad It will be observed that the grouping of the six tangents into pairs is arbitrary. By altering this grouping, we merely alter the linear function $l$, but do not obtain any new solution. Thus, say that the new form is $B = l'A + p'q'r'$, then, by properly determining the linear function $l$, we can reduce this to the original form $B = lA + pqr$; viz., we can satisfy identically the equation $(l - l')A + pqr - p'q'r'^2 = 0$; or what is the same thing, $\lambda A + pqr - p'q'r'^2 = 0$, where $\lambda$ is a linear function of the coordinates. We have, in fact, the conic $A = 0$ and the cubic $pqr = 0$ intersecting in the six points of contact; any other cubic through these six points; and consequently the cubic $p'q'r' = 0$ must be expressible in the form $\lambda A + pqr = 0$, and we have thus the identity in question.

\textbf{\S 138.}\quad We have just seen that the value of $B$ is necessarily of the form $B = lA + pqr$, but we are not concerned with its expression in this particular form. What we require in the sequel is a value of $B$, and thence one of $C$, satisfying the identity in question, $AC - B^2 = t_1t_2t_3t_4t_5t_6$; or what is the same thing, introducing for convenience a constant factor $K$, the identity
\[
AC - B^2 = Kt_1t_2t_3t_4t_5t_6.
\]

\textbf{\S 139.}\quad Instead of $C$, I write $\Gamma$, and consider the sextic amphigram $A\Gamma - B^2 = 0$, touched by the conic $A = 0$ in the points of contact of the conic with the six tangents $t_1 = 0$, $t_2 = 0$, $\ldots$, $t_6 = 0$. Suppose the sextic curve breaks up into factors; if one of these factors is a line, it is one of the six tangents, say the tangent $t_6 = 0$. If there is a conic factor, this is a conic touching the conic at its points of contact with two of the tangents, say the equation is $lA + mt_1t_2 = 0$. Similarly, if there is a cubic, quartic, or quintic factor, then the equation hereof is $sA + mt_1t_2t_3 = 0$, $UA + mt_1t_2t_3t_4 = 0$, or $VA + mt_1t_2t_3t_4t_5 = 0$. Or going on to the next case of a sextic factor (being of course the whole curve), we may say that this is $WA + mt_1t_2t_3t_4t_5t_6 = 0$. (Observe that since $AC - B^2 = Kt_1t_2t_3t_4t_5t_6$, this means only that the equation of the sextic amphigram is of the assumed form $A\Gamma - B^2 = 0$.)

\textbf{\S 140.}\quad By what precedes we can, for a sextic amphigram which breaks up in any assigned manner, determine the value of $\Gamma$. For instance, let the amphigram break up into two cubic curves; say these are $sA + mt_1t_2t_3 = 0$, $s'A + m't_4t_5t_6 = 0$. Assume
\[
A\Gamma - B^2 = mm'(sA + mt_1t_2t_3)(s'A + m't_4t_5t_6),
\]
then this equation is
\[
A\Gamma - B^2 = mm'\{ss'A^2 + (sm't_4t_5t_6 + s'mt_1t_2t_3)A\} + AC - B^2;
\]
that is, we have
\[
\Gamma = C + mm'(ss'A + sm't_4t_5t_6 + s'mt_1t_2t_3),
\]
and so in any other case.

I have already adverted to the question of the ``nodal determination'' of the formulae, and it might be properly here considered; viz., the question is as to the determination of the constants in such manner that, for instance, $sA + mt_1t_2t_3 = 0$ may be a nodal cubic, $UA + mt_1t_2t_3t_4 = 0$ a nodal, binodal, or trinodal quartic, \&c.; but I defer it for the moment in order first to apply the theory to the quartic surfaces.

\subsection*{Application to Quartic Surfaces.}

\textbf{\S 141.}\quad If a quartic surface has a node or nodes, we may take for a node the point $x = 0$, $y = 0$, $z = 0$; the equation of the quartic surface is then of the form
\[
Aw^2 + 2Bw + \Gamma = 0,
\]
where $A$, $B$, $\Gamma$ are functions of $x$, $y$, $z$ of the degrees 2, 3, 4 respectively. $A = 0$ is the tangent quadricone at the node in question; and the circumscribed cone is $A\Gamma - B^2 = 0$. By what precedes, this is an amphigram touching the quadricone along six generating lines thereof; say the tangent planes of the cone $A = 0$ along these six lines respectively are $t_1 = 0$, $t_2 = 0$, $\ldots$, $t_6 = 0$. We have then an identical equation
\[
AC - B^2 = Kt_1t_2t_3t_4t_5t_6,
\]
viz., regarding for a moment this equation as an equation for the determination of $B$, and $B'$ as any particular solution thereof, then its general solution is $B' + tA$, where $t$ is an arbitrary linear function of $(x, y, z)$, and the $B$ in the equation of the surface is properly $= B' + tA$. But by the substituting $w - t$ in place of $w$, the $B$ of the equation of the surface would then be made $= B'$; and it thus appears that we may, without loss of generality, take the $B$ of this equation to be any particular value satisfying the identity in question; and then, $B$ having such particular value, $C$ is a quartic function of $(x, y, z)$ completely determined by the same identity. And we then, by what precedes, at once determine $\Gamma$ so that the circumscribed cone $A\Gamma - B^2 = 0$ may be a cone breaking up in any assigned manner; for instance, if it be a cone $(3, 3)$, then, as just mentioned, the two cubic cones are $sA + mt_1t_2t_3 = 0$, $s'A + m't_4t_5t_6 = 0$; and $\Gamma$ has the value
\[
\Gamma = C + mm'(ss'A + sm't_4t_5t_6 + s'mt_1t_2t_3),
\]
above obtained.

\subsection*{On the Nodal Determination.}

\textbf{\S 142.}\quad I am not able to discuss with much completeness the question of nodal determination. We have to consider a cubic curve $sA + mt_1t_2t_3 = 0$, a quartic curve $UA + mt_1t_2t_3t_4 = 0$, \&c., as the case may be, and to determine the constants so that this shall have a node or nodes. Consider for a moment the form $PA + t_iQ = 0$, where $Q$ denotes the product $mt_1t_2\ldots$ of all or any of the tangents $t_1$, $\ldots$, $t_6$; the orders of $PA$, $t_iQ$ are of course equal, that is, the order of $P$ is less by unity than that of $Q$. I say that, by establishing a single relation between the constants, this may be made to have a node at the point of contact $A = 0$, $t_i = 0$. In fact, writing $A = \lambda\partial_\alpha + \mu\partial_\beta + \nu\partial_\gamma$, where $\lambda$, $\mu$, $\nu$ are arbitrary, there will be a node at any point if for that point $\frac{1}{\Delta}(PA + t_iQ) = 0$. But for the point $A = 0$, $t_i = 0$ this becomes $PA\Delta + Q\Delta t_i = 0$; moreover, if $t = 0$ be any other tangent of the conic $A = 0$, and if $p = 0$ be the line joining the points of contact of the tangents $t$, $t_i$, then we may write $A = tt_i - p^2$, and thence (since at the point in question, $A = 0$, $t_i = 0$, we have also $p = 0$) we find $\Delta tt_i = t\Delta t_i$, and the foregoing equation thus becomes $(tP + Q)t\Delta t_i = 0$; viz., this equation is satisfied irrespectively of the values of $\lambda$, $\nu$, if only at the point in question (that is, for the values of the coordinates which belong to the point $t_i = 0$, $A = 0$) we have $tP + Q = 0$, which is a single relation between the constants.

\textbf{\S 143.}\quad In particular the cubic curve $sA + mt_1t_2t_3 = 0$ may be made to have a node at the point of contact of any one of the three tangents; the quartic curve $UA + mt_1t_2t_3t_4 = 0$, a node, or two or three nodes, at the point or points of contact of any one, two, or three of the four tangents; and so in other cases. These are not the only solutions, and they are in fact solutions which (as afterwards explained) I propose to reject, attending in each case only to the remaining or proper solutions of the problem.

\textbf{\S 144.}\quad To obtain in a different manner the foregoing result, consider again the cubic curve $sA + mt_1t_2t_3 = 0$; regarding this as a given curve, the conic $A = 0$ is a conic determined (not of course completely) as a conic having therewith 3 points of 2-pointic intersection; viz., if the cubic has a node, then the cone $A = 0$ is either a conic passing through the node and besides touching the curve twice, or else it is a conic touching the curve 3 times; the former is of course the above mentioned case where there is a node at one of the points of contact on the conic $A = 0$; the latter is regarded as the proper solution. So in the case of a quartic curve $UA + mt_1t_2t_3t_4 = 0$, regarding this as a given curve, the conic $A = 0$ is a conic having therewith four points of 2-pointic intersection; viz., if the quartic curve has one, two, or three nodes, then the conic is either a conic passing through one, two, or three nodes, and besides touching the quartic thrice, twice, or once; or else it is a conic touching the quartic four times. The former is the above mentioned case where there is a node or nodes at a point or points of contact with the conic $A = 0$; the latter is regarded as the proper solution.

\textbf{\S 145.}\quad To fix the ideas, and at the same time obtain a result which will be afterwards useful, I work out the formulae for the cubic curve $sA + mt_1t_2t_3 = 0$, taking this equation under the form
\[
\left(\frac{\lambda}{x} + \frac{\mu}{y} + \frac{\nu}{z}\right)(x^2 + y^2 + z^2 - 2yz - 2zx - 2xy) + mxyz = 0.
\]
This may have a node in two different ways; viz.,

$1^{\circ}$. At the point of contact of one of the tangents $x = 0$, $y = 0$, $z = 0$ with the conic $A = 0$; say at the point of contact of $x = 0$, that is, the point $x = 0$, $y - z = 0$. The value of $m$ is $p = -\frac{4}{\mu}$; hence $\frac{1}{p} = -\frac{1}{4}\mu$; and, substituting, the equation of the curve becomes
\[
\left(\frac{\lambda}{x} + \frac{\mu}{y} + \frac{\nu}{z}\right)\{x^2 - 2x(y + z) + (y - z)^2\} + \frac{1}{4}\mu(x + y - z)^2 = 0,
\]
which has obviously a node at the point in question.

$2^{\circ}$. The node may be at a point not on the conic $A = 0$, viz. the value of $m$ is $= -\frac{4(\lambda + \mu + \nu)}{\lambda\mu\nu}$, the equation is
\[
\left(\frac{\lambda}{x} + \frac{\mu}{y} + \frac{\nu}{z}\right)(x^2 + y^2 + z^2 - 2yz - 2zx - 2xy) - \frac{4(\lambda + \mu + \nu)}{\lambda\mu\nu}xyz = 0.
\]
In fact, writing for shortness
\begin{align*}
\lambda + \mu + \nu &= L, & \lambda - \mu + \nu &= M, & \lambda + \mu - \nu &= N, \\
\lambda + \mu + \nu &= P, & -\lambda + \mu + \nu &= Q, & \&\text{c.},
\end{align*}
the node is at the point $x : y : z = L\lambda : M\mu : N\nu$; which is at once verified, if we remark that, writing for convenience $x, y, z = L\lambda, M\mu, N\nu$, then we have
\begin{align*}
-x + y + z &= MN, & x - y + z &= NL, & x + y - z &= LM, \\
\frac{\lambda}{x} + \frac{\mu}{y} + \frac{\nu}{z} &= P, \\
x^2 + y^2 + z^2 - 2yz - 2zx - 2xy &= -LMNP\ (= A).
\end{align*}
For, of the three equations for the coordinates of a node, the first is
\[
\frac{2\lambda}{x}A + P(-2y - 2z + 2x) - \frac{4(\lambda + \mu + \nu)}{\mu\nu}x = 0,
\]
that is, for the values in question,
\[
-\frac{2}{L}LMNP + P(-2MN) + \frac{4}{\mu\nu}LMN\lambda\nu = 0,
\]
that is
\[
-2LMNP - 2LMNP + 4LMNP = 0,
\]
which is satisfied; and similarly the other two equations are satisfied.

\subsection*{Quartic Surfaces resumed.}

\textbf{\S 146.}\quad Passing now from curves to cones, and to the theory of the quartic surface, suppose that there is a component cone having a nodal line, say the cubic cone $sA + mt_1t_2t_3 = 0$: if the remaining factor is $t_4t_5t_6$, then we have
\[
A\Gamma - B^2 = mm'(sA + mt_1t_2t_3)t_4t_5t_6.
\]
Suppose the nodal line is a line of contact with the cone $A = 0$, say its equations are $t_i = 0$, $p = 0$ ($p$ a linear function), then $sA + mt_1t_2t_3$ is a quadric function $(\mid\mid t_i, p)^2$, (of course with variable coefficients); hence $A\Gamma - B^2$ is a quadric function; and $A$ being a linear function $(\mid\mid t_i, p)$, it follows that $B$ is a linear function, and thence that $\Gamma$ is also a linear function; that is, $A$, $B$, $\Gamma$ are each of them a linear function $(\mid\mid t_i, p)$, or the line in question (viz.\ the line of contact $A = 0$, $t_i = 0$) is a line on the quartic surface $Aw^2 + 2Bw + \Gamma = 0$. As already mentioned, I exclude from consideration the surfaces which have upon them a line through two nodes; that is, I exclude from consideration the case in question where any component cone, or say where the sextic cone, has a nodal line which is a line on the tangent cone $A = 0$.

\textbf{\S 147.}\quad Now, excluding the case just referred to, I assume as a postulate that there is but one way in which the cubic cone $sA + mt_1t_2t_3 = 0$ can be made to have a nodal line, or the quartic cone $UA + mt_1t_2t_3t_4 = 0$ one, two, or three nodal lines, \&c., as the case may be. It is to be understood that this does not mean that the constants are in any of these cases completely determined, but that there is between them a relation or relations constituting a general solution which includes in itself every particular solution whatever. I have no doubt that as regards the cubic cone at least the assumption is correct. This being so, the character of a single node determines the nature of the surface; for instance, if there is a node $(3_1, 3)$, then taking this as the point $(x = 0,\; y = 0,\; z = 0)$ the equation of the surface is $Aw^2 + 2Bw + \Gamma = 0$, where
\[
\Gamma = C + mm'(ss'A + l'smt_4t_5t_6 + ls'mt_1t_2t_3),
\]
a surface of a determinate nature; so that the character of all the remaining nodes is completely determined.

\textbf{\S 148.}\quad The point to be attended to is, that if for instance there were two essentially distinct ways of giving the cubic cone $sA + mt_1t_2t_3 = 0$ a nodal line (such as there would be if the excluded case were considered admissible), then the foregoing equation of the surface would or might include two distinct forms of equation applying to different kinds of surface. The conclusion is that there is but one kind of quartic surface having a node $(3_1, 3)$. Admitting this, and similarly that there is but one kind of quartic surface having a node $6_{10}$, it follows that if (as the fact is) there is a surface having the nodes $1(6_{10}) + 10(3_1, 3)$ (Kummer's 11-nodal surface), then that the two first-mentioned kinds are in fact each of them this last-mentioned kind of surface; and it was in this manner that I arrived at the enumeration given near the beginning of the present Memoir.

\textbf{\S 149.}\quad The reasoning is, of course, in place of a direct demonstration which would consist in showing that a surface having a node $(3_1, 3)$ has 9 other like nodes, and also a node $6_{10}$; and that a surface having a node $6_{10}$ has 10 other nodes $(3_1, 3)$; and that, starting from either form of equation, we could, by passing to a node of the other kind, obtain the other form of equation.

\subsection*{Enumeration of the Cases.}

\textbf{\S 150.}\quad I collect the results as follows: I call to mind that we have always the identical equation $AC - B^2 = Kt_1t_2t_3t_4t_5t_6$, that the equation of the surface is $Aw^2 + 2Bw + \Gamma = 0$, and that the circumscribed cone is $A\Gamma - B^2 = 0$. The equation of a surface having different kinds of nodes will assume different forms according as the origin (or point $x = 0$, $y = 0$, $z = 0$) is taken to be at a node of one or other of these kinds; these forms of the equations are distinguished as ``node-forms,'' viz., we speak of the node-form $(3_1, 3)$ when the origin is a node $(3_1, 3)$, and so in other cases.

\vspace{0.5em}
\textbf{The 16-nodal surface}
\[16(1, 1, 1, 1, 1, 1),\]
\textbf{node-form} $(1, 1, 1, 1, 1, 1)$, cone is $6$,
and viz., equation is $Aw^2 + 2Bw + C = 0$.

\vspace{0.5em}
\textbf{The 15-nodal surface}
\[15(2, 1, 1, 1, 1),\]
\textbf{node-form} $(2, 1, 1, 1, 1)$, cone is $(lA + mt_1t_2t_3t_4t_5)t_6 = 0$,
and $\Gamma = C + \frac{K}{l}t_1t_2t_3t_4t_5t_6 = 0$.

\vspace{0.5em}
\textbf{The 14-nodal surface}
\[8(3_1, 1, 1, 1) + 6(2, 2, 1, 1),\]
\textbf{node-form} $(3_1, 1, 1, 1)$, cone is $(sA + mt_1t_2t_3)t_4t_5t_6 = 0$,
where $sA + mt_1t_2t_3 = 0$ is a nodal cubic $3_1$,
and $\Gamma = C + \frac{K}{m}t_4t_5t_6$;

\textbf{node-form} $(2, 2, 1, 1)$,
cone is $(lA + mt_1t_2)(l'A + m't_3t_4)t_5t_6 = 0$,
and $\Gamma = C + mm'(ll'A + lm't_3t_4 + l'mt_1t_2)t_5t_6$.

\vspace{0.5em}
\textbf{The 13($\alpha$)-nodal surface}
\[3(4_1, 1, 1) + 1(3, 1, 1, 1) + 9(3_1, 2, 1),\]
\textbf{node-form} $(4_1, 1, 1)$,
cone is $(UA + mt_1t_2t_3t_4)t_5t_6 = 0$,
where $UA + mt_1t_2t_3t_4 = 0$ is a trinodal quartic $4_1$,
and $\Gamma = C + \frac{K}{m}t_5t_6$;

\textbf{node-form} $(3, 1, 1, 1)$,
cone is $(sA + mt_1t_2t_3)t_4t_5t_6 = 0$,
and $\Gamma = C + \frac{K}{m}t_4t_5t_6$;

\textbf{node-form} $(3_1, 2, 1)$,
cone is $(sA + mt_1t_2t_3)(l'A + m't_4t_5)t_6 = 0$,
where $sA + mt_1t_2t_3 = 0$ is a nodal cubic $3_1$,
and $\Gamma = C + mm'(sl'A + m'st_4t_5t_6 + l'mt_1t_2t_3)t_6$.

\vspace{0.5em}
\textbf{The 13($\beta$)-nodal surface}
\[13(2, 2, 2),\]
\textbf{node-form} $(2, 2, 2)$,
cone is $(lA + mt_1t_2)(l'A + m't_3t_4)(l''A + m''t_5t_6) = 0$,
and
\begin{multline*}
\Gamma = C + mm'm''\{ll'l''A^2 + (ll'm''t_5t_6 + l''lm't_3t_4 + l'm''t_1t_2)A \\
+ lm'm''t_3t_4t_5t_6 + l'lm''t_1t_2t_5t_6 + l''m''t_1t_2t_3t_4\}.
\end{multline*}

\vspace{0.5em}
\textbf{The 12($\alpha$)-nodal surface}
\[12(4_1, 2),\]
\textbf{node-form} $(4_1, 2)$,
cone is $(UA + mt_1t_2t_3t_4)(l'A + m't_5t_6) = 0$,
where $UA + mt_1t_2t_3t_4 = 0$ is a trinodal quartic $4_1$,
and $\Gamma = C + mm'(l'UA + m'mt_3t_4t_5t_6 + l'mt_1t_2t_3t_4)$.

\vspace{0.5em}
\textbf{The 12($\beta$)-nodal surface}
\[2(5_1, 1) + 6(3_1, 3_1) + 4(3, 2, 1),\]
\textbf{node-form} $(5_1, 1)$,
cone is $(VA + mt_1t_2t_3t_4t_5)t_6 = 0$,
where $VA + mt_1t_2t_3t_4t_5 = 0$ is a 5-nodal quintic $5_1$,
and $\Gamma = C + \frac{K}{m}Vt_6$;

\textbf{node-form} $(3_1, 3_1)$,
cone is $(sA + mt_1t_2t_3)(s'A + m't_4t_5t_6) = 0$,
where $sA + mt_1t_2t_3 = 0$ and $s'A + m't_4t_5t_6 = 0$ are each of them a nodal cubic $3_1$,
and $\Gamma = C + mm'(ss'A + m'st_4t_5t_6 + ms't_1t_2t_3)$;

\textbf{node-form} $(3, 2, 1)$,
cone is $(sA + mt_1t_2t_3)(l'A + m't_4t_5)t_6 = 0$,
and $\Gamma = C + mm'(sl'A + m'st_4t_5t_6 + l'mt_1t_2t_3)t_6$.

\vspace{0.5em}
\textbf{The 12($\gamma$)-nodal surface,}
\[12(4_1, 1, 1),\]
\textbf{node-form} $(4_1, 1, 1)$,
cone is $(UA + mt_1t_2t_3t_4)t_5t_6 = 0$,
where $UA + mt_1t_2t_3t_4 = 0$ is a binodal quartic $4_1$,
and $\Gamma = C + \frac{K}{m}t_5t_6$.

\vspace{0.5em}
\textbf{The 11($\alpha$)-nodal surface,}
\[1(6_{10}) + 10(3_1, 3)\]
\textbf{node-form} $(6_{10})$,
cone is $WA + mt_1t_2t_3t_4t_5t_6 = 0$,
where this is a 10-nodal sextic $6_{10}$,
and $\Gamma = C + \frac{K}{m}W$;

\textbf{node-form} $(3_1, 3)$,
cone is $(sA + mt_1t_2t_3)(s'A + m't_4t_5t_6) = 0$,
where $sA + mt_1t_2t_3 = 0$ is a nodal cubic $3_1$,
and $\Gamma = C + mm'(ss'A + m'st_4t_5t_6 + ms't_1t_2t_3)$.

\vspace{0.5em}
\textbf{Other 11-nodal surfaces,}

\textbf{node-form} $(5_1, 1)$,
cone is $(VA + mt_1t_2t_3t_4t_5)t_6 = 0$,
where $VA + mt_1t_2t_3t_4t_5 = 0$ is a 5-nodal quintic $5_1$,
and $\Gamma = C + \frac{K}{m}Vt_6$;

\textbf{node-form} $(4_1, 2)$,
cone is $(UA + mt_1t_2t_3t_4)(l'A + m't_5t_6) = 0$,
where $UA + mt_1t_2t_3t_4 = 0$ is a binodal quartic $4_1$,
and $\Gamma = C + mm'(l'UA + m'Ut_3t_4 + l'mt_1t_2t_3t_4)$;

\textbf{node-form} $(4_1, 1, 1)$,
cone is $(UA + mt_1t_2t_3t_4)t_5t_6 = 0$,
where $UA + mt_1t_2t_3t_4$ is a nodal quartic $4_1$,
and $\Gamma = C + \frac{K}{m}t_5t_6$;

but whether these node-forms belong to the same or to different surfaces is not ascertained.

The enumeration is not extended to the 10-nodal surfaces, but I consider one case of these surfaces.

\vspace{0.5em}
\textbf{The 10($\alpha$)-nodal surface} $10(3, 3)$.

\textbf{\S 151.}\quad I assume only that there is a single node $(3, 3)$: taking the cone to be
\[
(sA + mt_1t_2t_3)(s'A + m't_4t_5t_6) = 0;
\]
then for the equation of the surface, in the node-form $(3, 3)$ in question, we have
\[
\Gamma = C + \frac{K}{mm'}(ss'A + sm't_4t_5t_6 + s'mt_1t_2t_3).
\]
But I present this result under a different form, as follows: I write
\[
A = p^2 - ft_1t_2 = q^2 - gt_3t_4 = r^2 - ht_5t_6,
\]
where $f$, $g$, $h$ are constants, and, as before, $p = 0$, $q = 0$, $r = 0$ are the lines joining the points of contact of $t_1$, $t_2$; $t_3$, $t_4$; and $t_5$, $t_6$ respectively: we have
\[
sA + mt_1t_2t_3 = sA + mt_3\left(A - \frac{p^2}{f}\right), \quad\text{and}\quad s'A + m't_4t_5t_6 = s'A + m't_6\left(A - \frac{r^2}{h}\right),
\]
or in place of $s$, $s'$ introducing new linear functions $\sigma$, $\sigma'$, the cubic curves may be taken to be $\sigma A - \frac{m}{f}p^2t_3$, $\sigma'A - \frac{m'}{h}r^2t_6$, so that we have
\begin{multline*}
A\Gamma - B^2 = mm'\left(\sigma A - \frac{m}{f}p^2t_3\right)\left(\sigma'A - \frac{m'}{h}r^2t_6\right) \\
= mm'\left\{\sigma\sigma'A^2 - \left(\frac{\sigma'm}{f}p^2t_3 + \frac{\sigma m'}{h}r^2t_6\right)A + \frac{mm'}{fh}p^2r^2t_3t_6\right\};
\end{multline*}
whence $B = \sqrt{mm'}\left(pqr + tA\right)$, where $t$ is a linear function of the coordinates; and we then have
\[
\Gamma = C + mm'\left(\sigma\sigma'A - \frac{\sigma'm}{f}p^2t_3 - \frac{\sigma m'}{h}r^2t_6\right) + \ldots
\]

\vspace{2em}
\begin{center}
{\small [To be continued.]}
\end{center}

\end{document}
